%
% This file is part of ncku-thesis-templete.
%
% ncku-thesis-templete is distributed in the hope that it will be useful,
% you can redistribute it and/or modify
% it under the terms of the Attribution-NonCommercial-ShareAlike
% 4.0 International.
%
% You should have received a copy of the
% Attribution-NonCommercial-ShareAlike 4.0 International
% along with ncku-thesis-templete.
%
% If not, see <http://creativecommons.org/licenses/by-nc-sa/4.0/legalcode.txt>.
%

% ------------------------------------------------

\documentclass[12pt, a4paper]{report}

% ------------------------------------------------

% XeLaTex 的設定

% --------------------------
% 加這個就可以設定字體
\usepackage{fontspec}

% 使用xeCJK，其他的還有CJK或是xCJK
\usepackage[BoldFont]{xeCJK}

% Set english font (Use package <ifplatform> to check)
%\usepackage{ifplatform}
%\ifwindows
  % In Windows
  \setmainfont{Times New Roman}
%\else
  % In *nix
%  \usepackage{times}
%  \usepackage[T1]{fontenc}
%\fi 

% 設定中英文的字型
% 字型的設定可以使用系統內的字型，而不用像以前一樣另外安裝
\setCJKmainfont{標楷體}
\newfontlanguage{Chinese}{CHN}
\setCJKfamilyfont{vchifont}[Script=CJK,Language=Chinese,Vertical=RotatedGlyphs]{標楷體}
\setCJKfamilyfont{normalchifont}[Script=CJK,Language=Chinese]{標楷體}

% 中文自動換行
\XeTeXlinebreaklocale "zh"

% 文字的彈性間距
\XeTeXlinebreakskip = 0pt plus 1pt

% 設定段落之間的距離
\setlength{\parskip}{0.3cm}

% 設定行距
\linespread{1.5}\selectfont

% Set style
\pagestyle{plain}

% ------------------------------------------------

% LaTex packages

% --------------------------

\usepackage{float}
\usepackage{pgfkeys}
\usepackage{graphicx}
\usepackage{hyperref}
\usepackage{pdfpages}
\usepackage{url}
\usepackage{multirow}
\newcommand{\tabincell}[2]{\begin{tabular}{@{}#1@{}}#2\end{tabular}}
\usepackage{fp}
\usepackage{siunitx}
\usepackage{ifthen}
\usepackage{xparse}
\usepackage{readarray}
\usepackage{framed}
\usepackage{caption}
\usepackage{subcaption}
\usepackage{amsmath}
\usepackage{booktabs}

\usepackage{comment}

\usepackage{algorithm}
\usepackage[noend]{algpseudocode}
\algnewcommand\algorithmicswitch{\textbf{switch}}
\algnewcommand\algorithmiccase{\textbf{case}}
\algnewcommand\algorithmicdefault{\textbf{default}}
\algnewcommand\algorithmicbreak{\textbf{break}}
\algdef{SE}[SWITCH]{Switch}{EndSwitch}[1]{\algorithmicswitch\ #1\ }{\algorithmicend\ \algorithmicswitch}%
\algdef{SE}[CASE]{Case}{EndCase}[1]{\algorithmiccase\ #1:}{\algorithmicend\ \algorithmiccase}%
\algdef{SE}[CASE]{Default}{EndDefault}[0]{\algorithmicdefault:}{\algorithmicend\ \algorithmiccase}%
\algtext*{EndSwitch}%
\algtext*{EndCase}%
\algtext*{EndDefault}%
\def\Break{\algorithmicbreak}

% Use in spine
\def\CJKmovesymbol#1{\raise.35em\hbox{#1}}
\def\CJKmove{\let\CJKsymbol\CJKmovesymbol
             \let\CJKpunctsymbol\CJKmovesymbol}

% ------------------------------------------------

% 有關學校對論文要求的設定

% --------------------------

% 學校排版 Arrangement style
%\usepackage{./ncku/style/macros}
\usepackage{./ncku/style/utdiss}

% --------------------------

% Sytle of itemize
\renewcommand{\labelitemi}{$\bullet$}
\renewcommand{\labelitemii}{$-$}
\renewcommand{\labelitemiii}{$\diamond$}
\renewcommand{\labelitemiv}{$\ast$}

% --------------------------

% 一些用來設定function和variable的command
%
% This file is part of ncku-thesis-templete.
%
% ncku-thesis-templete is distributed in the hope that it will be useful,
% you can redistribute it and/or modify
% it under the terms of the Attribution-NonCommercial-ShareAlike
% 4.0 International.
%
% You should have received a copy of the
% Attribution-NonCommercial-ShareAlike 4.0 International
% along with ncku-thesis-templete.
%
% If not, see <http://creativecommons.org/licenses/by-nc-sa/4.0/legalcode.txt>.
%

% ----------------------------------------------------------------------------
% 一些用來設定function和variable的command
% Some function and variable that let user use and configure
%
% 此處只是一些預設值和function
% 修改內容是在'conf/conf'
% ----------------------------------------------------------------------------

% Static variable and some provided API
%
% This file is part of ncku-thesis-templete.
%
% ncku-thesis-templete is distributed in the hope that it will be useful,
% you can redistribute it and/or modify
% it under the terms of the Attribution-NonCommercial-ShareAlike
% 4.0 International.
%
% You should have received a copy of the
% Attribution-NonCommercial-ShareAlike 4.0 International
% along with ncku-thesis-templete.
%
% If not, see <http://creativecommons.org/licenses/by-nc-sa/4.0/legalcode.txt>.
%

% Some common helper function

% ----------------------------------------------------------------------------

% Some helper functions

\newcommand{\GetMonthInEng}[1]
{
  \ifthenelse{\equal{#1}{1}}{January}{}
  \ifthenelse{\equal{#1}{2}}{February}{}
  \ifthenelse{\equal{#1}{3}}{March}{}
  \ifthenelse{\equal{#1}{4}}{April}{}
  \ifthenelse{\equal{#1}{5}}{May}{}
  \ifthenelse{\equal{#1}{6}}{June}{}
  \ifthenelse{\equal{#1}{7}}{July}{}
  \ifthenelse{\equal{#1}{8}}{August}{}
  \ifthenelse{\equal{#1}{9}}{September}{}
  \ifthenelse{\equal{#1}{10}}{October}{}
  \ifthenelse{\equal{#1}{11}}{November}{}
  \ifthenelse{\equal{#1}{12}}{December}{}
} % End of \newcommand{}

% 計算出台灣民國幾年
% Get the year using Taiwans' year
\newcommand{\SetOralTaiwanYear}[1]
{
  \FPeval{\OralTaiwanYearResult}{clip(#1 - 1911)}
} % End of \newcommand{}

% 計算出台灣民國幾年
% Get the year using Taiwans' year
\newcommand{\SetThesisTaiwanYear}[1]
{
  \FPeval{\ThesisTaiwanYearResult}{clip(#1 - 1911)}
} % End of \newcommand{}

% ----------------------------------------------------------------------------

% --- 關鍵字 Keyword ---
\newcommand{\KeywordsA}{} % Default
\newcommand{\KeywordsB}{} % Default
\newcommand{\KeywordsC}{} % Default
\newcommand{\KeywordsD}{} % Default
\newcommand{\KeywordsE}{} % Default
\newcommand{\SetKeywordsA}[1]{\renewcommand{\KeywordsA}{#1}}
\newcommand{\SetKeywordsB}[1]{\renewcommand{\KeywordsB}{#1}}
\newcommand{\SetKeywordsC}[1]{\renewcommand{\KeywordsC}{#1}}
\newcommand{\SetKeywordsD}[1]{\renewcommand{\KeywordsD}{#1}}
\newcommand{\SetKeywordsE}[1]{\renewcommand{\KeywordsE}{#1}}
\newcommand{\GetKeywordsA}[0]{\KeywordsA}
\newcommand{\GetKeywordsB}[0]{\KeywordsB}
\newcommand{\GetKeywordsC}[0]{\KeywordsC}
\newcommand{\GetKeywordsD}[0]{\KeywordsD}
\newcommand{\GetKeywordsE}[0]{\KeywordsE}
\DeclareDocumentCommand{\SetKeywords}{
  m G{\empty} G{\empty} G{\empty} G{\empty}}
{
  \SetKeywordsA{#1}
  \ifthenelse{\equal{#2}{\empty}}{}{\SetKeywordsB{#2}}
  \ifthenelse{\equal{#3}{\empty}}{}{\SetKeywordsC{#3}}
  \ifthenelse{\equal{#4}{\empty}}{}{\SetKeywordsD{#4}}
  \ifthenelse{\equal{#5}{\empty}}{}{\SetKeywordsE{#5}}
} % End of \newcommand{}

% ----------------------------------------------------------------------------

% In the minimal example below the macro \modulo{<a>}{<b>} stores the result of <a> mod <b> in the macro \result
\newcommand{\modulo}[2]{%
  \FPeval{\result}{trunc(#1-(#2*trunc(#1/#2,0)),0)}%
}

% ----------------------------------------------------------------------------



% Some helper function about insert image
\usepackage{ifthen}
\usepackage{xparse}

% ----------------------------------------------------------------------------
%
% http://tex.stackexchange.com/questions/34312/how-to-create-a-command-with-key-values
%
% 用\begin{figure} .. \end{figure}
% 可能會出現問題
% http://www.tex.ac.uk/cgi-bin/texfaq2html?label=ouparmd
%
% ----------------------------------------------------------------------------

\DeclareDocumentCommand{\SetImageCaption}{G{\empty}}
{
  \ifthenelse{\equal{#1}{\empty}}
  {}
  {\IfNoValueF{#1}{\caption{#1}}}
} % End of \DeclareDocumentCommand{}

\DeclareDocumentCommand{\SetImageLabel}{G{\empty}}
{
  \ifthenelse{\equal{#1}{\empty}}
  {}
  {\IfNoValueF{#1}{\label{#1}}}
} % End of \DeclareDocumentCommand{}

% -----------------------------------------------------------------

\pgfkeys
{
  /InsertImage/.is family, /InsertImage,
  default/.style = 
  {
    scale = 1.0,
    angle = 0,
    caption = \empty,
    label = \empty,
    align = \empty,
    pos = H,
    col = 1,
  },
  scale/.estore in = \TmpValueScale,
  angle/.estore in = \TmpValueAngle,
  caption/.estore in = \TmpValueCaption,
  label/.estore in = \TmpValueLabel,
  align/.estore in = \TmpValueAlign,
  pos/.estore in = \TmpValuePosition,
  col/.estore in = \TmpValueColumn,
} % End of \pgfkeys{}

% Insert a single column image
\newcommand{\InsertImage}[2][\empty]
{
  % Parse the input
  \pgfkeys{/InsertImage, default, #1}
  %
  % Print begin block
  \ifthenelse{\equal{\TmpValueColumn}{1}}
  { %
    \begin{figure}[H]
  } % End of if{}
  { %
    \begin{figure*}[Ht]
  } % End of else{}
  %
  % Do center if needed
  \ifthenelse{\equal{\TmpValueAlign}{center}} %
    {\center}{}
  %
  % Insert image
  \includegraphics %
    [scale=\TmpValueScale,%
      angle=\TmpValueAngle]{#2}%
  %
  % Set Caption
  \SetImageCaption{\TmpValueCaption}
  %
  % Set Label
  \ifthenelse{\equal{\TmpValueCaption}{\empty}} %
    {}{\SetImageLabel{\TmpValueLabel}}
  %
  % Print end block
  \ifthenelse{\equal{\TmpValueColumn}{1}}
  { %
    \end{figure}
  } % End of if{}
  { %
    \end{figure*}
  } % End of else{}
} % End of \newcommand{}

% -----------------------------------------------------------------

% ------------ OLD API ------------
% A wrapper to handle \InsertCenterImage{}
\newcommand{\InsertCenterImage}[2][\empty]
{
  \InsertImage[#1, align=center]{#2}
} % End of \newcommand{}

% -----------------------------------------------------------------


% Some helper function about insert image
\usepackage{ifthen}
\usepackage{xparse}

% ----------------------------------------------------------------------------
%
% http://tex.stackexchange.com/questions/34312/how-to-create-a-command-with-key-values
%
% 用\begin{figure} .. \end{figure}
% 可能會出現問題
% http://www.tex.ac.uk/cgi-bin/texfaq2html?label=ouparmd
%
% ----------------------------------------------------------------------------

% -----------------------------------------------------------------

% For multi images
\newcommand\MultiImagesPerRow{1}
\newcommand{\SetMultiImagesPerRow}[1]
  {\renewcommand{\MultiImagesPerRow}{#1}}
\newcommand{\GetMultiImagesPerRow}{\MultiImagesPerRow}

\newcommand\MultiImageTotalValue{0}
\newcommand{\SetMultiImageTotalValue}[1]
  {\renewcommand{\MultiImageTotalValue}{#1}}
\newcommand{\GetMultiImageTotalValue}{\MultiImageTotalValue}

\newcommand\MultiImageId{0}
\newcommand{\SetMultiImageId}[1]
  {\renewcommand{\MultiImageId}{#1}}
\newcommand{\GetMultiImageId}{\MultiImageId}

\newcommand\WidthOfImagePerRow{0}

%add desired spacing between images, e. g. ~, \quad, \qquad etc.
%(or a blank line to force the subfigure onto a new line)
\newcommand{\SubfigureBreakSpaceLine}[1]
{
  \ifthenelse{\equal{\intcalcMod{
    \GetMultiImageId - 1}{\GetMultiImagesPerRow}}{0}}%
  {%
    % Echo blank line
    %

  } % End of if{}
  {%
    ~
  } % End of else{}
} % End of \newcommand{}

\newcommand{\SetWidthOfImagePerRow}
{
  %
%  GetMultiImagesPerRow: \GetMultiImagesPerRow\\
%  GetMultiImageTotalValue: \GetMultiImageTotalValue\\
  %
  \ifthenelse{\equal{\GetMultiImagesPerRow}{1}}
  {
    %perrow = 1
    \FPeval{\WidthOfImagePerRow}{1.0}
  } % End of if{}
  {
    %perrow > 1
    %
    % Flag
    \provideboolean{WidthOfImagePerRowTmpIsDone}
    \setboolean{WidthOfImagePerRowTmpIsDone}{false}
    %
    % Calc remain
    \FPeval{\TmpRemain}{
      clip((\GetMultiImageTotalValue) - (\GetMultiImageId))}
    %TmpRemain: \TmpRemain\\
    %
    % Get mod
    \modulo{\GetMultiImageTotalValue}{\GetMultiImagesPerRow}
    %mod: \result\\
    %
    \FPiflt{\TmpRemain}{\GetMultiImagesPerRow}
    \else
      % Gt
      \setboolean{WidthOfImagePerRowTmpIsDone}{true}
      %
      % Default
      \FPeval{
        \WidthOfImagePerRow}{
          clip(clip(1 / (\GetMultiImagesPerRow)) - 0.1)}
    \fi % End of \FPiflt
    %
    \ifthenelse{\boolean{WidthOfImagePerRowTmpIsDone}}
    {}
    {
      % Lt
      \ifthenelse{\result = 0}
      {
        % Last row
        %
        \FPeval{
          \WidthOfImagePerRow}{
            clip(clip(1 / (\GetMultiImagesPerRow)) - 0.1)}
      } % End of if{}
      {
        % Diff row
        %
        \FPifeq{\TmpRemain}{0}
          \FPeval{\WidthOfImagePerRow}{1.0}
        \else
          \FPeval{
            \WidthOfImagePerRow}{
              clip(clip(1 / (\GetMultiImagesPerRow)) - 0.1)}
        \fi % End of \FPiflt
      } % End of else{}
    } % End of else{}
  } % End of else{}
  %
%  \WidthOfImagePerRow
  %
} % End of \newcommand{}

\pgfkeys
{
  /InsertMultiImages/.is family, /InsertMultiImages,
  default/.style = 
  {
    perrow = 1,
    caption = \empty,
    label = \empty,
    align = \empty,
    pos = H,
    col = 1,
  },
  perrow/.estore in = \TmpMIValueImagePerRow,
  caption/.estore in = \TmpMIValueCaption,
  label/.estore in = \TmpMIValueLabel,
  align/.estore in = \TmpMIValueAlign,
  pos/.estore in = \TmpMIValuePosition,
  col/.estore in = \TmpMIValueColumn,
} % End of \pgfkeys{}

% Insert multi-image
% Arg: 1st: Table configure
%      2~9th: Images (Max 8 Images)
%
% ---------------------------------------
% 使用 \InsertMultiImageLL#2
% 是跟 \InsertMultiImageLL{#2} 不一樣的
% 直接連接#2是指把#2整個當成function的所有參數
% 而{#2}是指把#2當成單一個參數傳給function
% ---------------------------------------
%
\DeclareDocumentCommand{\InsertMultiImages}{
  +O{\empty} +m +G{\empty} +G{\empty} +G{\empty}
  +G{\empty} +G{\empty} +G{\empty} +G{\empty}}
{
  % Parse the input
  \pgfkeys{/InsertMultiImages, default, #1}
  %
  %Set Multi Image Total
  \ifthenelse{\equal{#2}{\empty}}{}{\SetMultiImageTotalValue{1}}
  \ifthenelse{\equal{#3}{\empty}}{}{\SetMultiImageTotalValue{2}}
  \ifthenelse{\equal{#4}{\empty}}{}{\SetMultiImageTotalValue{3}}
  \ifthenelse{\equal{#5}{\empty}}{}{\SetMultiImageTotalValue{4}}
  \ifthenelse{\equal{#6}{\empty}}{}{\SetMultiImageTotalValue{5}}
  \ifthenelse{\equal{#7}{\empty}}{}{\SetMultiImageTotalValue{6}}
  \ifthenelse{\equal{#8}{\empty}}{}{\SetMultiImageTotalValue{7}}
  \ifthenelse{\equal{#9}{\empty}}{}{\SetMultiImageTotalValue{8}}
  %
  % Set per row
  \FPeval{\MultiImagesPerRow}{clip(
    \TmpMIValueImagePerRow)} %
  %
  %
  % Print begin block
  \ifthenelse{\equal{\TmpMIValueColumn}{1}}
  { %
    \begin{figure}[H]
  } % End of if{}
  { %
    \begin{figure*}[Ht]
  } % End of else{}
  %
  % Do center if needed
  \ifthenelse{\equal{\TmpMIValueAlign}{center}} %
    {\center}{}
    %
    % Image 1
    %\#2: P #2 P\\
    \ifthenelse{\equal{#2}{\empty}}
      {}
      {
        \SetMultiImageId{1}
        \SubfigureBreakSpaceLine{1}
        \InsertSubfigure#2
      } % End of else{}
    %
    % Image 2
    %\#3: P #3 P\\
    \ifthenelse{\equal{#3}{\empty}}
      {}
      {
        \SetMultiImageId{2}
        \SubfigureBreakSpaceLine{2}
        \InsertSubfigure#3
      } % End of else{}
    %
    % Image 3
    %\#4: P #4 P\\
    \ifthenelse{\equal{#4}{\empty}}
      {}
      {
        \SetMultiImageId{3}
        \SubfigureBreakSpaceLine{3}
        \InsertSubfigure#4
      } % End of else{}
    %
    % Image 4
    %\#5: P #5 P\\
    \ifthenelse{\equal{#5}{\empty}}
      {}
      {
        \SetMultiImageId{4}
        \SubfigureBreakSpaceLine{4}
        \InsertSubfigure#5
      } % End of else{}
    %
    % Image 5
    %\#6: P #6 P\\
    \ifthenelse{\equal{#6}{\empty}}
      {}
      {
        \SetMultiImageId{5}
        \SubfigureBreakSpaceLine{5}
        \InsertSubfigure#6
      } % End of else{}
    %
    % Image 6
    %\#7: P #7 P\\
    \ifthenelse{\equal{#7}{\empty}}
      {}
      {
        \SetMultiImageId{6}
        \SubfigureBreakSpaceLine{6}
        \InsertSubfigure#7
      } % End of else{}
    %
    % Image 7
    %\#8: P #8 P\\
    \ifthenelse{\equal{#8}{\empty}}
      {}
      {
        \SetMultiImageId{7}
        \SubfigureBreakSpaceLine{7}
        \InsertSubfigure#8
      } % End of else{}
    %
    % Image 8
    %\#9: P #9 P\\
    \ifthenelse{\equal{#9}{\empty}}
      {}
      {
        \SetMultiImageId{8}
        \SubfigureBreakSpaceLine{8}
        \InsertSubfigure#9
      } % End of else{}
    %
    % Set Caption
    \SetImageCaption{\TmpMIValueCaption}
    % Set Label
    \ifthenelse{\equal{\TmpMIValueCaption}{\empty}}
      {}{\SetImageLabel{\TmpMIValueLabel}}
    %
  % Print end block
  \ifthenelse{\equal{\TmpMIValueColumn}{1}}
  { %
    \end{figure}
  } % End of if{}
  { %
    \end{figure*}
  } % End of else{}
} % End of \newcommand{}

\pgfkeys
{
  /InsertSubfigure/.is family, /InsertSubfigure,
  default/.style = 
  {
    scale = 1.0,
    angle = 0,
    caption = \empty,
    label = \empty,
    align = \empty,
  },
  scale/.estore in = \TmpMISubValueScale,
  angle/.estore in = \TmpMISubValueAngle,
  caption/.estore in = \TmpMISubValueCaption,
  label/.estore in = \TmpMISubValueLabel,
  align/.estore in = \TmpMISubValueAlign,
} % End of \pgfkeys{}

% Low-level insert image
\newcommand{\InsertSubfigure}[2][\empty]
{
  % Set width
  \SetWidthOfImagePerRow
%  \WidthOfImagePerRow
  %
  % Parse the input
  \pgfkeys{/InsertSubfigure, default, #1}
  %
  \ifthenelse{\equal{\TmpMIValueColumn}{2}}
  { %
    \begin{subfigure}{\WidthOfImagePerRow\textwidth}
  } % End of if{}
  { %
    \begin{subfigure}{\WidthOfImagePerRow\columnwidth}
  } % End of else{}
    %
  %
  % Do center if needed
  \ifthenelse{\equal{\TmpMISubValueAlign}{center}} %
    {\center}{}
    %
    \includegraphics[
      scale=\TmpMISubValueScale,
      angle=\TmpMISubValueAngle]{#2}
    %
    % Set Caption
    \SetImageCaption{\TmpMISubValueCaption}
    % Set Label
    \ifthenelse{\equal{
      \TmpMISubValueCaption}{\empty}}
      {}{\SetImageLabel{\TmpMISubValueLabel}}
  \end{subfigure}%
} % End of \newcommand{}

% ----------------------------------------------------------------------------

%
% This file is part of ncku-thesis-templete.
%
% ncku-thesis-templete is distributed in the hope that it will be useful,
% you can redistribute it and/or modify
% it under the terms of the Attribution-NonCommercial-ShareAlike
% 4.0 International.
%
% You should have received a copy of the
% Attribution-NonCommercial-ShareAlike 4.0 International
% along with ncku-thesis-templete.
%
% If not, see <http://creativecommons.org/licenses/by-nc-sa/4.0/legalcode.txt>.
%

% Some helper function use for oral document

% ----------------------------------------------------------------------------
\def \OralIndexHeader {Oral presentation document}
% ----------------------------------------------------------------------------

% 口試委員 Committee member(s)
\newcommand{\CommitteeSize}{8} % Default
\newcommand{\GetCommitteeSize}[0]{\CommitteeSize}
\newcommand{\SetCommitteeSize}[1]
{
  \ifthenelse{#1 < 2}
  {
    \renewcommand{\CommitteeSize}{2}
  } % End of if{}
  {
    \ifthenelse{#1 > 8}
    {\renewcommand{\CommitteeSize}{8}}
    {\renewcommand{\CommitteeSize}{#1}}
  } % End of else{}
} % End of \newcommand{}

% 口試委員簽名區
\newcommand{\DisplayCommitteeSignatureArea}
{
  % 口試委員 至少2位
  \ifthenelse{\CommitteeSize = 2}
  {
    \begin{minipage}[c][6.4cm][c]{\textwidth}
      \makebox[0.5\textwidth][c]{\namesigdate}
      \makebox[0.5\textwidth][c]{\namesigdate}
    \end{minipage}
  } % End of if{}
  {} % End of else{}

  \ifthenelse{\CommitteeSize = 3}
  {
    \begin{minipage}[c][3.2cm][c]{\textwidth}
      \makebox[0.5\textwidth][c]{\namesigdate}
      \makebox[0.5\textwidth][c]{\namesigdate}
    \end{minipage}

    \begin{minipage}[c][3.2cm][c]{\textwidth}
      \makebox[\textwidth][c]{\namesigdate}
    \end{minipage}
  } % End of if{}
  {} % End of else{}

  \ifthenelse{\CommitteeSize = 4}
  {
    \begin{minipage}[c][3.2cm][c]{\textwidth}
      \makebox[0.5\textwidth][c]{\namesigdate}
      \makebox[0.5\textwidth][c]{\namesigdate}
    \end{minipage}

    \begin{minipage}[c][3.2cm][c]{\textwidth}
      \makebox[0.5\textwidth][c]{\namesigdate}
      \makebox[0.5\textwidth][c]{\namesigdate}
    \end{minipage}
  } % End of if{}
  {} % End of else{}

  \ifthenelse{\CommitteeSize = 5}
  {
    \begin{minipage}[c][2.13cm][c]{\textwidth}
      \makebox[0.5\textwidth][c]{\namesigdate}
      \makebox[0.5\textwidth][c]{\namesigdate}
    \end{minipage}

    \begin{minipage}[c][2.13cm][c]{\textwidth}
      \makebox[0.5\textwidth][c]{\namesigdate}
      \makebox[0.5\textwidth][c]{\namesigdate}
    \end{minipage}

    \begin{minipage}[c][2.13cm][c]{\textwidth}
      \makebox[\textwidth][c]{\namesigdate}
    \end{minipage}
  } % End of if{}
  {} % End of else{}

  \ifthenelse{\CommitteeSize = 6}
  {
    \begin{minipage}[c][2.13cm][c]{\textwidth}
      \makebox[0.5\textwidth][c]{\namesigdate}
      \makebox[0.5\textwidth][c]{\namesigdate}
    \end{minipage}

    \begin{minipage}[c][2.13cm][c]{\textwidth}
      \makebox[0.5\textwidth][c]{\namesigdate}
      \makebox[0.5\textwidth][c]{\namesigdate}
    \end{minipage}

    \begin{minipage}[c][2.13cm][c]{\textwidth}
      \makebox[0.5\textwidth][c]{\namesigdate}
      \makebox[0.5\textwidth][c]{\namesigdate}
    \end{minipage}
  } % End of if{}
  {} % End of else{}

  \ifthenelse{\CommitteeSize = 7}
  {
    \begin{minipage}[c][1.6cm][c]{\textwidth}
      \makebox[0.5\textwidth][c]{\namesigdate}
      \makebox[0.5\textwidth][c]{\namesigdate}
    \end{minipage}

    \begin{minipage}[c][1.6cm][c]{\textwidth}
      \makebox[0.5\textwidth][c]{\namesigdate}
      \makebox[0.5\textwidth][c]{\namesigdate}
    \end{minipage}

    \begin{minipage}[c][1.6cm][c]{\textwidth}
      \makebox[0.5\textwidth][c]{\namesigdate}
      \makebox[0.5\textwidth][c]{\namesigdate}
    \end{minipage}

    \begin{minipage}[c][1.6cm][c]{\textwidth}
      \makebox[\textwidth][c]{\namesigdate}
    \end{minipage}
  } % End of if{}
  {} % End of else{}

  \ifthenelse{\CommitteeSize = 8}
  {
    \begin{minipage}[c][1.6cm][c]{\textwidth}
      \makebox[0.5\textwidth][c]{\namesigdate}
      \makebox[0.5\textwidth][c]{\namesigdate}
    \end{minipage}

    \begin{minipage}[c][1.6cm][c]{\textwidth}
      \makebox[0.5\textwidth][c]{\namesigdate}
      \makebox[0.5\textwidth][c]{\namesigdate}
    \end{minipage}

    \begin{minipage}[c][1.6cm][c]{\textwidth}
      \makebox[0.5\textwidth][c]{\namesigdate}
      \makebox[0.5\textwidth][c]{\namesigdate}
    \end{minipage}

    \begin{minipage}[c][1.6cm][c]{\textwidth}
      \makebox[0.5\textwidth][c]{\namesigdate}
      \makebox[0.5\textwidth][c]{\namesigdate}
    \end{minipage}
  } % End of if{}
  {} % End of else{}
} % End of \newcommand{}


% ----------------------------------------------------------------------------

% Signature line
\newcommand{\namesigdate}[0]{\rule{5.5cm}{1pt}}

% ----------------------------------------------------------------------------

% 口試証明文件 Document of oral presentation

% 要顯示圖片還是範例
\newcommand{\OralDocumentImage}{0}
\newcommand{\OralDocumentTemplete}{1}
\newcommand{\OralDocumentEngOnly}{0}
\newcommand{\OralDocument}{\OralDocumentTemplete} % Default

\newcommand{\DisplayOralTemplete}[0]
  {\renewcommand{\OralDocument}{\OralDocumentTemplete}}
\newcommand{\DisplayOralImage}[0]
  {\renewcommand{\OralDocument}{\OralDocumentImage}}
\newcommand{\DisplayOralEngOnly}[0]
  {\renewcommand{\OralDocumentEngOnly}{1}}

% The path of the image that the oral document
\newcommand\OralDocumentImageChiPath{./example/oral/oral-chi.pdf} % Default
\newcommand\OralDocumentImageEngPath{./example/oral/oral-eng.pdf} % Default
\newcommand{\SetOralImageChi}[1]
{\renewcommand{\OralDocumentImageChiPath}{./context/oral/#1}}
\newcommand{\SetOralImageEng}[1]
{\renewcommand{\OralDocumentImageEngPath}{./context/oral/#1}}

\newcommand{\SetOralImage}[2]
{
  \SetOralImageChi{#1}
  \SetOralImageEng{#2}
} % End of \newcommand{}

\newcommand{\GetOralImageChiPath}[0]{\OralDocumentImageChiPath}
\newcommand{\GetOralImageEngPath}[0]{\OralDocumentImageEngPath}

% ----------------------------------------------------------------------------

% Use to include the oral files
\newcommand{\DisplayOral}{%
% This file is part of ncku-thesis-templete.
%
% ncku-thesis-templete is distributed in the hope that it will be useful,
% you can redistribute it and/or modify
% it under the terms of the Attribution-NonCommercial-ShareAlike
% 4.0 International.
%
% You should have received a copy of the
% Attribution-NonCommercial-ShareAlike 4.0 International
% along with ncku-thesis-templete.
%
% If not, see <http://creativecommons.org/licenses/by-nc-sa/4.0/legalcode.txt>.
%

% ----------------------------------------------------------------------------
%                           Document of oral presentation
%                                   口試証明文件
% ----------------------------------------------------------------------------

% ----------------------------------------------------------------------------

\begin{comment}
\ifthenelse{\OralDocument = \OralDocumentTemplete}
{
  % -------------------------- 顯示範本 --------------------------
  \ifthenelse{\LangEnableChi = 1}
  {
    \StartNewPage

    % 設定顯示在目錄
    % If combine with english version
    \ifthenelse{\GetLangEnableEng = 1}
    {
      \addcontentsline{toc}{chapter}{Document of oral presentation}
    }{}

    % Chinese version
    \input{./ncku/oral/templete/templete-chi}

    % ------------------------------------------------
    \ifthenelse{\GetLangEnableDemo = 1}
    {
      \input{./ncku/oral/templete/templete-chi-2}
      \input{./ncku/oral/templete/templete-chi-3}
      \input{./ncku/oral/templete/templete-chi-4}
    }{}
    % ------------------------------------------------
  }{}

  % English page start
  \ifthenelse{\GetLangEnableEng = 1}
  {
    \StartNewPage

    % 設定顯示在目錄
    % If combine with chinese version
    \ifthenelse{\LangEnableChi = 1}
    {
      \addcontentsline{toc}{section}{English version}
    } % End of if{}
    {
      \addcontentsline{toc}{chapter}{Document of oral presentation}
    } % End of else{}

    \thispagestyle{empty}
    \input{./ncku/oral/templete/templete-eng}
  }{}
} % End of if{}
{
  % ---------------------- 顯示圖片, Image in PDF format ----------------------

  % Page start
  \StartNewPage
  \thispagestyle{empty}

  % 設定顯示在目錄
  \addcontentsline{toc}{chapter}{Document of oral presentation}

  % Chinese page start
  \ifthenelse{\LangEnableChi = 1}
  {
    \addcontentsline{toc}{section}{Chinese version}
    \includepdf{\GetOralImageChiPath}
  }{}

  % English page start

  % If combine with chinese version
  \ifthenelse{\LangEnableChi = 1}
  {
    \StartNewPage
    \thispagestyle{empty}
    \addcontentsline{toc}{section}{English version}
  }{}

  \includepdf{\GetOralImageEngPath}
} % End of else{}
\end{comment}

% ----------------------------------------------------------------------------

\ifthenelse{\OralDocument = \OralDocumentTemplete}
{
  % -------------------------- 顯示範本 --------------------------

  % 設定顯示在目錄
  \StartNewPage
  %\thispagestyle{empty}
  \addcontentsline{toc}{chapter}{\OralIndexHeader}

  % Chinese version
  \input{./ncku/oral/templete/templete-chi-5}

  % English version
  \StartNewPage
  %\thispagestyle{empty}
  \addcontentsline{toc}{section}{English version}
  \input{./ncku/oral/templete/templete-eng-2}
} % End of if{}
{
  % ---------------------- 顯示圖片, Image in PDF format ----------------------

  % 設定顯示在目錄
  \StartNewPage
  \thispagestyle{empty}
  \addcontentsline{toc}{chapter}{\OralIndexHeader}

  % Chinese version
  \addcontentsline{toc}{section}{Chinese version}
  \includepdf{\GetOralImageChiPath}

  % English version
  \StartNewPage
  \thispagestyle{empty}
  \addcontentsline{toc}{section}{English version}
  \includepdf{\GetOralImageEngPath}
} % End of else{}

% ----------------------------------------------------------------------------

% End of page
\EndOfPage

% ----------------------------------------------------------------------------
}

% ----------------------------------------------------------------------------

%
% This file is part of ncku-thesis-templete.
%
% ncku-thesis-templete is distributed in the hope that it will be useful,
% you can redistribute it and/or modify
% it under the terms of the Attribution-NonCommercial-ShareAlike
% 4.0 International.
%
% You should have received a copy of the
% Attribution-NonCommercial-ShareAlike 4.0 International
% along with ncku-thesis-templete.
%
% If not, see <http://creativecommons.org/licenses/by-nc-sa/4.0/legalcode.txt>.
%

% Some common helper function

% ----------------------------------------------------------------------------

% 使用 hyperref 在 pdf 簡介欄裡填入相關資料
\newcommand{\FillInPDFData}[0]
{
  \ifx \hypersetup \undefined
    % do nothing
    \relax
  \else
    \ifx \GetChiTitle \undefined
      \hypersetup
      {
        pdftitle  = {\GetEngTitle},
        pdfauthor = {\GetAuthorEngName},
      }
    \else
      \hypersetup
      {
        pdftitle  = {\GetEngTitle \enspace (\GetChiTitle)},
        pdfauthor = {\GetAuthorEngName \enspace (\GetAuthorChiName)},
      }
    \fi

    \hypersetup
    {
      unicode     = true,
      pdfcreator  = {\GetSchoolEngName},
      pdfproducer = {\GetSchoolEngName},
      pdfsubject  = {},
    }

    \ifthenelse{\equal{\GetKeywordsE}{\empty}}
    {
      \ifthenelse{\equal{\GetKeywordsD}{\empty}}
      {
        \ifthenelse{\equal{\GetKeywordsC}{\empty}}
        {
          \ifthenelse{\equal{\GetKeywordsB}{\empty}}
          {
            \ifthenelse{\equal{\GetKeywordsA}{\empty}}
            {}
            {
              \hypersetup
              {pdfkeywords = {\GetKeywordsA}}
            } % End of else{}
          } % End of if{}
          {
            \hypersetup
            {pdfkeywords = {\GetKeywordsA, \GetKeywordsB}}
          } % End of else{}
        } % End of if{}
        {
          \hypersetup
          {pdfkeywords = {\GetKeywordsA, \GetKeywordsB, \GetKeywordsC}}
        } % End of else{}
      } % End of if{}
      {
        \hypersetup
        {pdfkeywords = {\GetKeywordsA, \GetKeywordsB, \GetKeywordsC, \GetKeywordsD}}
      } % End of else{}
    } % End of if{}
    {
      \hypersetup
      {pdfkeywords = {\GetKeywordsA, \GetKeywordsB, \GetKeywordsC, \GetKeywordsD, \GetKeywordsE}}
    } % End of else{}
  \fi
} % End of \newcommand{}

% ----------------------------------------------------------------------------


%
% This file is part of ncku-thesis-templete.
%
% ncku-thesis-templete is distributed in the hope that it will be useful,
% you can redistribute it and/or modify
% it under the terms of the Attribution-NonCommercial-ShareAlike
% 4.0 International.
%
% You should have received a copy of the
% Attribution-NonCommercial-ShareAlike 4.0 International
% along with ncku-thesis-templete.
%
% If not, see <http://creativecommons.org/licenses/by-nc-sa/4.0/legalcode.txt>.
%

% Some helper function about equation

% ----------------------------------------------------------------------------

\newcommand{\SetEquationLabel}[1]
{
  \SetImageLabel{#1}
} % End of \newcommand{}

\DeclareDocumentCommand{\EquationBegin}{G{\empty}}
{
  \begin{equation}
  \SetEquationLabel{#1}
  \begin{aligned}
} % End of \newcommand{}

\newcommand{\EquationEnd}
{
  \end{aligned}
  \end{equation}
} % End of \newcommand{}

% ----------------------------------------------------------------------------


%
% This file is part of ncku-thesis-templete.
%
% ncku-thesis-templete is distributed in the hope that it will be useful,
% you can redistribute it and/or modify
% it under the terms of the Attribution-NonCommercial-ShareAlike
% 4.0 International.
%
% You should have received a copy of the
% Attribution-NonCommercial-ShareAlike 4.0 International
% along with ncku-thesis-templete.
%
% If not, see <http://creativecommons.org/licenses/by-nc-sa/4.0/legalcode.txt>.
%

% Some helper function for reference

% ----------------------------------------------------------------------------

% For equation
\newcommand{\RefEquation}[1]
{
  \ref{#1}
} % End of \newcommand{}

% For equation
\newcommand{\RefEquationB}[1]
{
  \eqref{#1}
} % End of \newcommand{}

% For bib
\newcommand{\RefBib}[1]
{
  \cite{#1}
} % End of \newcommand{}

% For figure, table, chapter, section, subsection, .etc
\newcommand{\RefTo}[1]
{
  \ref{#1}
} % End of \newcommand{}

\newcommand{\RefFigure}[1]
{
  \RefTo{#1}
} % End of \newcommand{}

\newcommand{\RefTable}[1]
{
  \RefTo{#1}
} % End of \newcommand{}

% For page
\newcommand{\RefPage}[1]
{
  \pageref{#1}
} % End of \newcommand{}

% ----------------------------------------------------------------------------



% Helper function for different page or chapter
%
% This file is part of ncku-thesis-templete.
%
% ncku-thesis-templete is distributed in the hope that it will be useful,
% you can redistribute it and/or modify
% it under the terms of the Attribution-NonCommercial-ShareAlike
% 4.0 International.
%
% You should have received a copy of the
% Attribution-NonCommercial-ShareAlike 4.0 International
% along with ncku-thesis-templete.
%
% If not, see <http://creativecommons.org/licenses/by-nc-sa/4.0/legalcode.txt>.
%

% Some helper function about thesis

% ----------------------------------------------------------------------------

\newcommand{\BeginThesis}[0]
{
  % Start of paper
  \begin{document}

  % Set the line spacing for whole thesis
  \baselineskip = 26pt
} % End of \newcommand{}

\newcommand{\EndThesis}[0]
{
  % End of paper
  \end{document}
} % End of \newcommand{}

\begin{comment}
\newcommand{\IncludeContext}[0]
{
  % 預設使用模板的說明文件(./example/context)為內容
  % 請修改"./context/context"中以成為你的論文內容
  \ifthenelse{\GetLangEnableDemo = 1}%
  {% ------------------------------------------------

% 封面 Cover
\DisplayCover

% ------------------------------------------------

% 口試証明文件 Document of oral presentation
\DisplayOral

% ------------------------------------------------

% 摘要 Abstract
% 請根據你的需求去選用中文或英文版
%
% 學校要求如果是寫中文版的話
% 則要同時交出英文版
%
% 只交英文版就不用交中文版
%
% 詳細學校的要求可去附件看
\input{./context/abstract/abstract}

% ------------------------------------------------

% 致謝 Acknowledgments
\input{./context/acknowledgments/acknowledgments}

% ------------------------------------------------

% 目錄 (內容, 圖表和圖片) Index of contents, tables and figures.
% 內容會自動產生 The indices will generate in automate.
\DisplayIndex

% ------------------------------------------------

% Introduction chapter
\input{./context/chp1/chp1}

% Objective chapter
\input{./context/chp2/chp2}

% Related Work chapter
\input{./context/chp3/chp3}

% Algorithm chapter
%\input{./context/chp4/chp4}

% Performance chapter
%\input{./context/chp5/chp5}

% Conclusion chapter
%\input{./context/conclusion/conclusion}

% Future work chapter
%\input{./context/future-work/future-work}

% ------------------------------------------------

% 參考文獻 References
\input{./context/references/references}

% ------------------------------------------------

% 附錄 Appendix
\input{./context/appendix/appendix}

% ------------------------------------------------
}%
  {% ------------------------------------------------

% 封面 Cover
\DisplayCover

% ------------------------------------------------

% 口試証明文件 Document of oral presentation
\DisplayOral

% ------------------------------------------------

% 摘要 Abstract
% 請根據你的需求去選用中文或英文版
%
% 學校要求如果是寫中文版的話
% 則要同時交出英文版
%
% 只交英文版就不用交中文版
%
% 詳細學校的要求可去附件看
\input{./context/abstract/abstract}

% ------------------------------------------------

% 致謝 Acknowledgments
\input{./context/acknowledgments/acknowledgments}

% ------------------------------------------------

% 目錄 (內容, 圖表和圖片) Index of contents, tables and figures.
% 內容會自動產生 The indices will generate in automate.
\DisplayIndex

% ------------------------------------------------

% Introduction chapter
\input{./context/chp1/chp1}

% Objective chapter
\input{./context/chp2/chp2}

% Related Work chapter
\input{./context/chp3/chp3}

% Algorithm chapter
%\input{./context/chp4/chp4}

% Performance chapter
%\input{./context/chp5/chp5}

% Conclusion chapter
%\input{./context/conclusion/conclusion}

% Future work chapter
%\input{./context/future-work/future-work}

% ------------------------------------------------

% 參考文獻 References
\input{./context/references/references}

% ------------------------------------------------

% 附錄 Appendix
\input{./context/appendix/appendix}

% ------------------------------------------------
}
} % End of \newcommand{}
\end{comment}

\newcommand{\CreateThesis}[0]
{
  % Start of thesis
  \BeginThesis

  % 內容 context
  \ifthenelse{\GetLangEnableDemo = 1}%
  {% ------------------------------------------------

% 封面 Cover
\DisplayCover

% ------------------------------------------------

% 口試証明文件 Document of oral presentation
\DisplayOral

% ------------------------------------------------

% 摘要 Abstract
% 請根據你的需求去選用中文或英文版
%
% 學校要求如果是寫中文版的話
% 則要同時交出英文版
%
% 只交英文版就不用交中文版
%
% 詳細學校的要求可去附件看
\input{./context/abstract/abstract}

% ------------------------------------------------

% 致謝 Acknowledgments
\input{./context/acknowledgments/acknowledgments}

% ------------------------------------------------

% 目錄 (內容, 圖表和圖片) Index of contents, tables and figures.
% 內容會自動產生 The indices will generate in automate.
\DisplayIndex

% ------------------------------------------------

% Introduction chapter
\input{./context/chp1/chp1}

% Objective chapter
\input{./context/chp2/chp2}

% Related Work chapter
\input{./context/chp3/chp3}

% Algorithm chapter
%\input{./context/chp4/chp4}

% Performance chapter
%\input{./context/chp5/chp5}

% Conclusion chapter
%\input{./context/conclusion/conclusion}

% Future work chapter
%\input{./context/future-work/future-work}

% ------------------------------------------------

% 參考文獻 References
\input{./context/references/references}

% ------------------------------------------------

% 附錄 Appendix
\input{./context/appendix/appendix}

% ------------------------------------------------
}%
  {% ------------------------------------------------

% 封面 Cover
\DisplayCover

% ------------------------------------------------

% 口試証明文件 Document of oral presentation
\DisplayOral

% ------------------------------------------------

% 摘要 Abstract
% 請根據你的需求去選用中文或英文版
%
% 學校要求如果是寫中文版的話
% 則要同時交出英文版
%
% 只交英文版就不用交中文版
%
% 詳細學校的要求可去附件看
\input{./context/abstract/abstract}

% ------------------------------------------------

% 致謝 Acknowledgments
\input{./context/acknowledgments/acknowledgments}

% ------------------------------------------------

% 目錄 (內容, 圖表和圖片) Index of contents, tables and figures.
% 內容會自動產生 The indices will generate in automate.
\DisplayIndex

% ------------------------------------------------

% Introduction chapter
\input{./context/chp1/chp1}

% Objective chapter
\input{./context/chp2/chp2}

% Related Work chapter
\input{./context/chp3/chp3}

% Algorithm chapter
%\input{./context/chp4/chp4}

% Performance chapter
%\input{./context/chp5/chp5}

% Conclusion chapter
%\input{./context/conclusion/conclusion}

% Future work chapter
%\input{./context/future-work/future-work}

% ------------------------------------------------

% 參考文獻 References
\input{./context/references/references}

% ------------------------------------------------

% 附錄 Appendix
\input{./context/appendix/appendix}

% ------------------------------------------------
}

  % End of thesis
  \EndThesis
} % End of \newcommand{}

% ----------------------------------------------------------------------------


%
% This file is part of ncku-thesis-templete.
%
% ncku-thesis-templete is distributed in the hope that it will be useful,
% you can redistribute it and/or modify
% it under the terms of the Attribution-NonCommercial-ShareAlike
% 4.0 International.
%
% You should have received a copy of the
% Attribution-NonCommercial-ShareAlike 4.0 International
% along with ncku-thesis-templete.
%
% If not, see <http://creativecommons.org/licenses/by-nc-sa/4.0/legalcode.txt>.
%

% Some helper function about page

% ----------------------------------------------------------------------------

\newcommand{\StartNewPage}
{
  % Page start
  \newpage
  \phantomsection
} % End of \newcommand{}

\newcommand{\EndOfPage}
{
  % End of page
  \clearpage
} % End of \newcommand{}

% ----------------------------------------------------------------------------


%
% This file is part of ncku-thesis-templete.
%
% ncku-thesis-templete is distributed in the hope that it will be useful,
% you can redistribute it and/or modify
% it under the terms of the Attribution-NonCommercial-ShareAlike
% 4.0 International.
%
% You should have received a copy of the
% Attribution-NonCommercial-ShareAlike 4.0 International
% along with ncku-thesis-templete.
%
% If not, see <http://creativecommons.org/licenses/by-nc-sa/4.0/legalcode.txt>.
%

% Some helper function use in cover

% ----------------------------------------------------------------------------

% --- University name 學校名字 ---
% 基本上是寫死, 但是如果是別校的人, 可直接使用\SetSchoolName來修改
\newcommand\univCname{國立成功大學}           % Default
\newcommand\univDeptCname{東海大學高階經營管理碩士在職專班(研究所)}           % Default
\newcommand\univEname{National Cheng Kung University} % Default
\newcommand{\SetSchoolChiName}[1]{\renewcommand{\univCname}{#1}}
\newcommand{\SetSchoolEngName}[1]{\renewcommand{\univEname}{#1}}
\newcommand{\SetSchoolName}[2]
{
  \SetSchoolChiName{#1}
  \SetSchoolEngName{#2}
} % End of \newcommand{}

\newcommand{\GetSchoolChiName}[0]{\univCname}
\newcommand{\GetSchoolDeptChiName}[0]{\univDeptCname}
\newcommand{\GetSchoolEngName}[0]{\univEname}
% ----------------------------------------------------------------------------

% --- Chinese / English title 中英文論文題目 ---
\newcommand{\cTitle}{Chinese Title Here} % Default
\newcommand{\eTitle}{English Title Here} % Default
\newcommand{\SetChiTitle}[1]{\renewcommand{\cTitle}{#1}}
\newcommand{\SetEngTitle}[1]{\renewcommand{\eTitle}{#1}}
\newcommand{\SetTitle}[2]
{
  \SetChiTitle{#1}
  \SetEngTitle{#2}
} % End of \newcommand{}

\newcommand{\GetChiTitle}[0]{\cTitle}
\newcommand{\GetEngTitle}[0]{\eTitle}
% ----------------------------------------------------------------------------

% --- User's name 使用者名字 ---
\newcommand{\myCname}{你的名字}     % Default
\newcommand{\myEname}{Your name}   % Default
\newcommand{\SetMyChiName}[1]{\renewcommand{\myCname}{#1}}
\newcommand{\SetMyEngName}[1]{\renewcommand{\myEname}{#1}}
\newcommand{\SetMyName}[2]
{
  \SetMyChiName{#1}
  \SetMyEngName{#2}
} % End of \newcommand{}

\newcommand{\GetAuthorChiName}[0]{\myCname}
\newcommand{\GetAuthorEngName}[0]{\myEname}

% ----------------------------------------------------------------------------

% --- Degree name 學位 ---
% thesis 是指論文的通稱
% dissertation 指的是博士的論文

% 碩士論文是             Master's thesis
% 但是一般碩士論文都會說   Master's dissertation

% 博士論文也可以說是      Doctoral thesis
% 但是一般博士論文都會說   Doctoral dissertation

\newcommand{\degreeCname}{碩士/博士} % Default
\newcommand{\degreeEname}{Master / Doctor} % Default
\newcommand{\degreeThesisEname}{Master's / Doctoral dissertation} % Default

\newcommand{\SetChiDegree}[1]{\renewcommand{\degreeCname}{#1}}
\newcommand{\SetEngDegree}[1]{\renewcommand{\degreeEname}{#1}}
\newcommand{\SetEngDegreeThesis}[1]{\renewcommand{\degreeThesisEname}{#1}}

\newcommand{\PhdDegree}[0]
{
  \SetChiDegree{博士}
  \SetEngDegree{Doctor}
  \SetEngDegreeThesis{Doctoral dissertation}
} % End of \newcommand{}

\newcommand{\MasterDegree}[0]
{
  \SetChiDegree{碩士}
  \SetEngDegree{Master}
  \SetEngDegreeThesis{Master's dissertation}
} % End of \newcommand{}

\newcommand{\GetChiDegree}[0]{\degreeCname}
\newcommand{\GetEngDegree}[0]{\degreeEname}
\newcommand{\GetEngDegreeThesis}[0]{\degreeThesisEname}
% ----------------------------------------------------------------------------

% --- Date 日期 ---

% --- 論文的日期 ---
\newcommand{\ThesisYear}{2014}  % Default
\newcommand{\ThesisMonth}{1}    % Default

\newcommand{\SetThesisDate}[2]
{
  \SetThesisTaiwanYear{#1}
  \renewcommand{\ThesisYear}{#1}
  \renewcommand{\ThesisMonth}{#2}
} % End of \newcommand{}

\newcommand{\GetThesisYear}[0]{\ThesisYear}
\newcommand{\GetThesisYearInTaiwanYear}[0]{\ThesisTaiwanYearResult}
\newcommand{\GetThesisMonth}[0]{\ThesisMonth}
\newcommand{\GetThesisMonthInEng}[0]{\GetMonthInEng{\ThesisMonth}}

% ---  口試的日期 ---
\newcommand{\OralChiYear}{101}      % Default
\newcommand{\OralChiMonth}{1}       % Default
\newcommand{\OralChiDay}{1}         % Default
\newcommand{\OralEngYear}{2014}     % Default
\newcommand{\OralEngMonth}{January} % Default
\newcommand{\OralEngDay}{1}         % Default

\newcommand{\GetOralChiYear}[0]{\OralChiYear}
\newcommand{\GetOralChiMonth}[0]{\OralChiMonth}
\newcommand{\GetOralChiDay}[0]{\OralChiDay}
\newcommand{\GetOralEngYear}[0]{\OralEngYear}
\newcommand{\GetOralEngMonth}[0]{\OralEngMonth}
\newcommand{\GetOralEngDay}[0]{\OralEngDay}

\newcommand{\SetOralChiDate}[3]
{
  \SetOralTaiwanYear{#1}
  \renewcommand{\OralChiYear}{\OralTaiwanYearResult}
  \renewcommand{\OralChiMonth}{#2}
  \renewcommand{\OralChiDay}{#3}
} % End of \newcommand{}

\newcommand{\SetOralEngDate}[3]
{
  \renewcommand{\OralEngYear}{#1}
  \renewcommand{\OralEngMonth}{\GetMonthInEng{#2}}
  \renewcommand{\OralEngDay}{#3}
} % End of \newcommand{}

\newcommand{\SetOralDate}[3]
{
  \SetOralChiDate{#1}{#2}{#3}
  \SetOralEngDate{#1}{#2}{#3}
} % End of \newcommand{}

% ----------------------------------------------------------------------------

% --- 學院 College, 系所 Department and Institute ---

% --------------------------- College ---------------------------
\newcommand{\collCname}{學院 C}
\newcommand{\collEname}{College of C}
\newcommand{\SetCollChiName}[1]{\renewcommand{\collCname}{#1}}
\newcommand{\SetCollEngName}[1]{\renewcommand{\collEname}{#1}}
\newcommand{\SetCollName}[2]
{
  \SetCollChiName{#1}
  \SetCollEngName{#2}
} % End of \newcommand{}

\newcommand{\GetCollChiName}[0]{\collCname}
\newcommand{\GetCollEngName}[0]{\collEname}

% --------------------------- Department ---------------------------
\newcommand{\deptCname}{A 系 / 所}
%\newcommand{\deptEname}{DeptA} % Short form of department
\newcommand{\fulldeptEname}{Department / Insitute A} % Full name of department
\newcommand{\SetDeptChiName}[1]{\renewcommand{\deptCname}{#1}}
%\newcommand{\SetDeptEngShortName}[1]{\renewcommand{\deptEname}{#1}}
\newcommand{\SetDeptEngName}[1]{\renewcommand{\fulldeptEname}{#1}}
\newcommand{\SetDeptName}[3]
{
  \SetDeptChiName{#1}
%  \SetDeptEngShortName{#2}
  \SetDeptEngName{#3}
} % End of \newcommand{}

\newcommand{\GetDeptChiName}[0]{\deptCname}
\newcommand{\GetDeptEngName}[0]{\fulldeptEname}

% ----------------------------------------------------------------------------

% --- 指導老師 Advisor(s) ---
% 在封面上預算了最多3位的空間
% 英文名字固定以Prof.開頭

\newcommand{\advisorCnameA}{X 教授}
\newcommand{\advisorEnameA}{Professor X}
\newcommand{\advisorCnameB}{}
\newcommand{\advisorEnameB}{}
\newcommand{\advisorCnameC}{}
\newcommand{\advisorEnameC}{}

\newcommand{\GetAdvisorChiNameA}{\advisorCnameA}
\newcommand{\GetAdvisorEngNameA}{\advisorEnameA}
\newcommand{\GetAdvisorChiNameB}{\advisorCnameB}
\newcommand{\GetAdvisorEngNameB}{\advisorEnameB}
\newcommand{\GetAdvisorChiNameC}{\advisorCnameC}
\newcommand{\GetAdvisorEngNameC}{\advisorEnameC}

\newcommand{\SetAdvisorChiNameA}[1]{\renewcommand{\advisorCnameA}{#1}}
\newcommand{\SetAdvisorEngNameA}[1]{\renewcommand{\advisorEnameA}{#1}}
\newcommand{\SetAdvisorChiNameB}[1]{\renewcommand{\advisorCnameB}{#1}}
\newcommand{\SetAdvisorEngNameB}[1]{\renewcommand{\advisorEnameB}{#1}}
\newcommand{\SetAdvisorChiNameC}[1]{\renewcommand{\advisorCnameC}{#1}}
\newcommand{\SetAdvisorEngNameC}[1]{\renewcommand{\advisorEnameC}{#1}}

\newcommand{\SetAdvisorNameA}[2]
{
  \SetAdvisorChiNameA{#1}
  \SetAdvisorEngNameA{#2}
} % End of \newcommand{}

\newcommand{\SetAdvisorNameB}[2]
{
  \SetAdvisorChiNameB{#1}
  \SetAdvisorEngNameB{#2}
} % End of \newcommand{}

\newcommand{\SetAdvisorNameC}[2]
{
  \SetAdvisorChiNameC{#1}
  \SetAdvisorEngNameC{#2}
} % End of \newcommand{}

% ----------------------------------------------------------------------------

% Use to include the cover files
\newcommand{\DisplayCover}{%
% This file is part of ncku-thesis-templete.
%
% ncku-thesis-templete is distributed in the hope that it will be useful,
% you can redistribute it and/or modify
% it under the terms of the Attribution-NonCommercial-ShareAlike
% 4.0 International.
%
% You should have received a copy of the
% Attribution-NonCommercial-ShareAlike 4.0 International
% along with ncku-thesis-templete.
%
% If not, see <http://creativecommons.org/licenses/by-nc-sa/4.0/legalcode.txt>.
%

% ------------------------------------------------

% 根據user的需求去選用中文或英文封面

%\ifthenelse{\GetLangEnableChi = 1}
%{
%  \input{./ncku/cover/cover-chi}
  \input{./ncku/cover/cover-chi-2-ok}
  \input{./ncku/cover/cover-blank}
  \input{./ncku/cover/cover-rpt}
%}{}

%\ifthenelse{\GetLangEnableDemo = 1}
%{
%  \input{./ncku/cover/cover-chi-2}
%}{}

%\ifthenelse{\GetLangEnableEng = 1}
%{
%  \input{./ncku/cover/cover-eng}
%  \input{./ncku/cover/cover-eng-3}
%}{}

% ------------------------------------------------
}

% ----------------------------------------------------------------------------

% Display Chinese and English name in english cover
\newcommand{\CoverDisplayNameChiEng}{0} % Default not display both

\newcommand{\SetCDBothName}[0]{\renewcommand{\CoverDisplayNameChiEng}{1}}

\newcommand{\CDBothName}[0]
{
  \SetCDBothName
} % End of \newcommand{}

\newcommand{\GetCDBothName}[0]{\CoverDisplayNameChiEng}

% ----------------------------------------------------------------------------

%
% This file is part of ncku-thesis-templete.
%
% ncku-thesis-templete is distributed in the hope that it will be useful,
% you can redistribute it and/or modify
% it under the terms of the Attribution-NonCommercial-ShareAlike
% 4.0 International.
%
% You should have received a copy of the
% Attribution-NonCommercial-ShareAlike 4.0 International
% along with ncku-thesis-templete.
%
% If not, see <http://creativecommons.org/licenses/by-nc-sa/4.0/legalcode.txt>.
%

% Some helper function about Spine

% ----------------------------------------------------------------------------

% Set width of the Spine
%\SetSpineWidth{3.0}

% 書脊 Spine
% 5X頁 精裝本 1.3cm
% 5X頁 普通本 0.6cm
% 以25頁來當成通用可以比較好
% 所以應該 ~0.3cm
% (Default: 1.5cm, minmum: 1.5cm, maximum: 2.5cm)
\newcommand{\SpineWidth}{0.4} % Default
\newcommand{\SpineColWidth}{0.1} % Default
\newcommand{\SpineEngColWidth}{0.1} % Default
\newcommand{\GetSpineWidth}[0]{\SpineWidth}
\newcommand{\GetSpineColWidth}[0]{\SpineColWidth}
\newcommand{\GetSpineEngColWidth}[0]{\SpineEngColWidth}
\newcommand{\SetSpineWidth}[1]
{
  % Update the width (Maximum)
  \FPifgt{#1}{2.5}
    \FPeval{\SpineWidth}{clip(2.5)}
  \else
%    \renewcommand{\SpineWidth}{#1}
    \FPeval{\SpineWidth}{clip(#1)}
  \fi
  %
  % Update the width (Minimum)
  \FPiflt{#1}{1.0}
    \FPeval{\SpineWidth}{clip(1.0)}
  \else
%    \renewcommand{\SpineWidth}{#1}
    \FPeval{\SpineWidth}{clip(#1)}
  \fi
  %
  % Update the column width
%  \FPeval{\SpineColWidth}{clip((#1 * 2) / 3) + clip(#1 / 5)}
  \FPeval{\SpineColWidth}{
    clip(\GetSpineWidth / 2) + clip(\GetSpineWidth / 5)}
  \FPeval{\SpineEngColWidth}{
    clip(\GetSpineWidth / 2) + clip(\GetSpineWidth / 7)}
} % End of \newcommand{}

% Title
\newcommand{\SpineTitleUseChi}{1}
\newcommand{\SpineTitleUseEng}{0}
\newcommand{\SpineTitleLang}{\SpineTitleUseEng} % Default
\newcommand{\GetSpineTitleLang}[0]{\SpineTitleLang}
\newcommand{\SpineTitleChi}[0]
{
  \renewcommand{\SpineTitleLang}{\SpineTitleUseChi}
} % End of \newcommand{}

% ----------------------------------------------------------------------------

% Use to include the oral files
\newcommand{\CreateSpine}
{
  \BeginThesis

  % Use for chinese version
  \ifthenelse{\GetLangEnableChi = 1}
  {
    %
% This file is part of ncku-thesis-templete.
%
% ncku-thesis-templete is distributed in the hope that it will be useful,
% you can redistribute it and/or modify
% it under the terms of the Attribution-NonCommercial-ShareAlike
% 4.0 International.
%
% You should have received a copy of the
% Attribution-NonCommercial-ShareAlike 4.0 International
% along with ncku-thesis-templete.
%
% If not, see <http://creativecommons.org/licenses/by-nc-sa/4.0/legalcode.txt>.
%

% ----------------------------------------------------------------------------
%    The spine of the book
%    書脊 (Chinese version)
% ----------------------------------------------------------------------------

% 先把Logo關掉
\ClearWatermark

\StartNewPage

% 設定使用 無頁碼, 有浮水印
\thispagestyle{empty}

% ------------------------------------------------

\ifthenelse{\GetLangEnableChi = 1}
{
  \StartNewPage

  % 設定使用 無頁碼, 有浮水印
  \thispagestyle{empty}

  {
    \CJKfamily{vchifont}
    \fontsize{12}{12}\selectfont

    \fbox{
      \parbox[c][0.99\textheight][c]{\GetSpineWidth cm}
      {
        \begin{tabular}{c}
          \rotatebox{-90}{
            \parbox[c][\GetSpineColWidth cm][c]{{0.15\textheight}
            }{\vspace{0.25cm}
            \scriptsize \GetSchoolChiName \\ \GetDeptChiName
          }}
          \\
          \rotatebox{-90}{
            \parbox[c][\GetSpineColWidth cm][c]{{0.05\textheight}
            }{\vspace{0.25cm}
            \scriptsize \centering \makebox{\GetChiDegree 論文}
          }}
          \\
          {
              \CJKfamily{normalchifont}
              \rotatebox{-90}{
                \parbox[c][\GetSpineColWidth cm][c]{{0.5\textheight}
                }{\vspace{0.45cm}
                \centering \parbox{0.5\textheight}%
                {%
                  \ifthenelse{\SpineTitleLang = \SpineTitleUseChi}%
                  {%
                    \ifthenelse{\GetLangEnableDemo = 1}
                    {\mbox{\tiny \GetEngTitle}}%
                    {\tiny \GetChiTitle}%
                  }%
                  {\mbox{\tiny \GetEngTitle}}%
                }%
              }}
              \CJKfamily{vchifont}
          }
          \\
          \rotatebox{-90}{
            \parbox[c][\GetSpineColWidth cm][c]{{0.08\textheight}
            }{\vspace{0.25cm}
            \centering \makebox{\scriptsize
              \GetAuthorChiName \enspace 撰}
          }}
          \\
          \rotatebox{-90}{
            \parbox[c][\GetSpineColWidth cm][c]{{0.08\textheight}
            }{
            \vspace{0.25cm} \scriptsize \centering
            {
              \fontspec[RawFeature={vertical}]{Times New Roman}
              \GetThesisYearInTaiwanYear
            }
            學年度
          }}
        \end{tabular}
      } % End of \parbox{}
    } % End of \fbox{}

  } % End of {}

  \EndOfPage
}{} % End of \ifthenelse{}

% ------------------------------------------------

\EndOfPage

% ------------------------------------------------

% If this is debug mode, then show in chinese title version
\ifthenelse{\GetLangEnableDemo = 1}
{
  \StartNewPage

  % 設定使用 無頁碼, 有浮水印
  \thispagestyle{empty}

  {
    \CJKfamily{vchifont}
    \fontsize{12}{12}\selectfont

    \fbox{
      \parbox[c][0.99\textheight][c]{\GetSpineWidth cm}
      {
        \begin{tabular}{c}
          \rotatebox{-90}{
            \parbox[c][\GetSpineColWidth cm][c]{{0.15\textheight}
            }{\vspace{0.25cm}
            \scriptsize \GetSchoolChiName \\ \GetDeptChiName
          }}
          \\
          \rotatebox{-90}{
            \parbox[c][\GetSpineColWidth cm][c]{{0.05\textheight}
            }{\vspace{0.25cm}
            \scriptsize \centering \makebox{\GetChiDegree 論文}
          }}
          \\
          {
              \CJKfamily{normalchifont}
              \rotatebox{-90}{
                \parbox[c][\GetSpineColWidth cm][c]{{0.5\textheight}
                }{\vspace{0.45cm}
                \centering \parbox{0.5\textheight}{
                \tiny \GetChiTitle}
              }}
              \CJKfamily{vchifont}
          }
          \\
          \rotatebox{-90}{
            \parbox[c][\GetSpineColWidth cm][c]{{0.08\textheight}
            }{\vspace{0.25cm}
            \centering \makebox{\scriptsize
              \GetAuthorChiName \enspace 撰}
          }}
          \\
          \rotatebox{-90}{
            \parbox[c][\GetSpineColWidth cm][c]{{0.08\textheight}
            }{
            \vspace{0.25cm} \scriptsize \centering
            {
              \fontspec[RawFeature={vertical}]{Times New Roman}
              \GetThesisYearInTaiwanYear
            }
            學年度
          }}
        \end{tabular}
      } % End of \parbox{}
    } % End of \fbox{}

  } % End of {}

  \EndOfPage
}{} % End of \ifthenelse{}

% ------------------------------------------------

  }{}

  % Use for english version
  \ifthenelse{\GetLangEnableEng = 1}
  {
    %
% This file is part of ncku-thesis-templete.
%
% ncku-thesis-templete is distributed in the hope that it will be useful,
% you can redistribute it and/or modify
% it under the terms of the Attribution-NonCommercial-ShareAlike
% 4.0 International.
%
% You should have received a copy of the
% Attribution-NonCommercial-ShareAlike 4.0 International
% along with ncku-thesis-templete.
%
% If not, see <http://creativecommons.org/licenses/by-nc-sa/4.0/legalcode.txt>.
%

% ----------------------------------------------------------------------------
%    The spine of the book
%    書脊 (English version)
% ----------------------------------------------------------------------------

% 先把Logo關掉
\ClearWatermark

\StartNewPage

% 設定使用 無頁碼, 有浮水印
\thispagestyle{empty}

% ------------------------------------------------

{
  %\CJKfamily{vchifont}
  \fontsize{8}{8}\selectfont

  \fbox{
    \parbox[c][0.99\textheight][c]{\GetSpineWidth cm}
    {
      \begin{tabular}{c}
        \rotatebox{-90}{
          \parbox[c][\GetSpineEngColWidth cm][c]{{0.23\textheight}
          }{\vspace{0.3cm}
          \tiny \GetSchoolEngName \\ \GetDeptEngName
        }}
        \\
        \hspace{0.01cm}
        \rotatebox{-90}{
          \parbox[c][\GetSpineEngColWidth cm][c]{{0.08\textheight}
          }{\vspace{0.3cm}
          \centering \makebox{\tiny \GetEngDegreeThesis}
        }}
        \\
        \hspace{0.02cm}
        \rotatebox{-90}{
          \parbox[c][\GetSpineEngColWidth cm][c]{{0.46\textheight}
          }{\vspace{0.3cm}
          \centering \parbox{0.46\textheight}{\mbox{\tiny \GetEngTitle}}
        }}
        \\
        \hspace{0.01cm}
        \rotatebox{-90}{
          \parbox[c][\GetSpineEngColWidth cm][c]{{0.09\textheight}
          }{\vspace{0.3cm}
          \centering \makebox{\tiny \GetAuthorEngName}
        }}
        \\
        \hspace{0.01cm}
        \rotatebox{-90}{
          \parbox[c][\GetSpineEngColWidth cm][c]{{0.05\textheight}
          }{\vspace{0.3cm}
          \centering \tiny \GetThesisMonthInEng \\
            \GetThesisYear
        }}
      \end{tabular}
    } % End of \parbox{}
  } % End of \fbox{}
} % End of {}

% ------------------------------------------------

\EndOfPage

% ------------------------------------------------

  }{}

  \EndThesis
} % End of \newcommand{}

% ----------------------------------------------------------------------------

%
% This file is part of ncku-thesis-templete.
%
% ncku-thesis-templete is distributed in the hope that it will be useful,
% you can redistribute it and/or modify
% it under the terms of the Attribution-NonCommercial-ShareAlike
% 4.0 International.
%
% You should have received a copy of the
% Attribution-NonCommercial-ShareAlike 4.0 International
% along with ncku-thesis-templete.
%
% If not, see <http://creativecommons.org/licenses/by-nc-sa/4.0/legalcode.txt>.
%

% Some helper function about chapter and section

% ----------------------------------------------------------------------------

% Chapters
\DeclareDocumentCommand{\StartChapter}{+m +g}
{
  \chapter{#1}
  \IfNoValueF{#2}{\label{#2}}
%  \baselineskip=26pt
%  \thispagestyle{empty}
} % End of \DeclareDocumentCommand{}

\def \EndChapter {\EndOfPage}

% Section
\DeclareDocumentCommand{\StartSection}{+m +g}
{
  \section{#1}
  \IfNoValueF{#2}{\label{#2}}
} % End of \DeclareDocumentCommand{}

% Sub-Section
\DeclareDocumentCommand{\StartSubSection}{+m +g}
{
  \subsection{#1}
  \IfNoValueF{#2}{\label{#2}}
} % End of \DeclareDocumentCommand{}

% Sub-Sub-Section
\DeclareDocumentCommand{\StartSubSubSection}{+m +g}
{
  \subsubsection{#1}
  \IfNoValueF{#2}{\label{#2}}
} % End of \DeclareDocumentCommand{}

% ----------------------------------------------------------------------------


%
% This file is part of ncku-thesis-templete.
%
% ncku-thesis-templete is distributed in the hope that it will be useful,
% you can redistribute it and/or modify
% it under the terms of the Attribution-NonCommercial-ShareAlike
% 4.0 International.
%
% You should have received a copy of the
% Attribution-NonCommercial-ShareAlike 4.0 International
% along with ncku-thesis-templete.
%
% If not, see <http://creativecommons.org/licenses/by-nc-sa/4.0/legalcode.txt>.
%

% Some helper function use in abstract

% ----------------------------------------------------------------------------

% Abstract
\newcommand{\StartChiAbstract}
{
  \StartNewPage

  % Add to "Table of Contents"
  \addcontentsline{toc}{chapter}{中文摘要}

  % Title
  %\centerline{\Large \textbf{摘要}}
\begin{minipage}[c][5cm][t]{\textwidth}
  論文名稱：社群組織之重要構成因素探討:  成員分享行為及\\ 
  \indent \hspace{5.0em} 輔助分享資訊工具研究  ---以g0v零時政府零時政府為例 \\
  校所名稱：\GetSchoolDeptChiName \\
  畢業時間：2015年 6月 \\
  研究生  ： \GetAuthorChiName \hspace{12.0em} 指導教授： \GetAdvisorChiNameA \\
  
  %
  %\makebox[4.0em][r]{校所名稱：} %abstract
  %\makebox[6.0em][l]{ \GetSchoolDeptChiName } \\
  %\vspace{0.5cm}  
  %\hspace{2.0em}
  %\makebox[4.0em][r]{畢業時間：}
  %\makebox[6.0em][l]{2015年 6月 } \\
  \vspace{0.5cm}  
  %\hspace{2.0em}
  %\makebox[4.0em][r]{研究生：}
  %\makebox[6.0em][l]{\GetAuthorChiName }  
  \vspace{0.5cm}  
  \makebox[8.0em][c]{}
  %\makebox[4.0em][r]{指導教授：}
  %\makebox[8.0em][l]{\GetAdvisorChiNameA}
  
  \hspace{2.0em}
  
  
\end{minipage}  

  \leftline{\Large \textbf{中文摘要:}}
  
  
  
} % End of \newcommand{}

\newcommand{\StartAbstract}
{
  \StartNewPage

  % Add to "Table of Contents"
  \addcontentsline{toc}{chapter}{英文摘要}

  % Title
\begin{minipage}[c][5cm][t]{\textwidth}
  
  Title of Thesis  ：The most important factor of psudo community :Majoring in \\
  \indent \hspace{6.5em} sharing behavior and tools ---Let g0v as example \\
  Name of Institute：Tunghai University  Executive Master of \\
  \indent \hspace{6.5em} Business Administration Program \\
  Graduation Time  ：06/2015 \\
  Student Name     ： \GetAuthorChiName \hspace{12.0em} Advisor Name： \GetAdvisorChiNameA \\

\end{minipage} 
  
  \leftline{\Large \textbf{Abstract:}}
} % End of \newcommand{}

\newcommand{\EndChiAbstract}
{
  % Keyword
  \AbstractKeyword

  \EndOfPage
} % End of \newcommand{}

\newcommand{\EndAbstract}
{
  % Keyword
  \ExtendedAbstractKeyword

  \EndOfPage
} % End of \newcommand{}

% Extended Abstract
\newcommand{\StartExtendedAbstract}
{
  \StartNewPage

  % Set page
  \baselineskip=20pt
  \setlength{\parindent}{0.0pt}

  % Set style
  \pagestyle{empty}

  % 設定段落之間的距離
  \setlength{\parskip}{0.5cm}

  % Add to "Table of Contents"
  \addcontentsline{toc}{chapter}{Extended Abstract}

  % Add title
  \parbox{\textwidth}{\center \large \textbf{\eTitle}}
  \vspace{0.5cm}

  % Add name
  \centerline{\large \GetAuthorEngName}

  % Add names
  \centerline{Prof. \thinspace \GetAdvisorEngNameA}
  \ifthenelse{\equal{\GetAdvisorEngNameB}{\empty}}
    {}
    {\centerline{Prof. \thinspace \GetAdvisorEngNameB}}

  \ifthenelse{\equal{\GetAdvisorEngNameC}{\empty}}
    {}
    {\centerline{Prof. \thinspace \GetAdvisorEngNameC}}

  % Add department
  \centerline{\GetDeptEngName}
  \centerline{\GetCollEngName}
} % End of \newcommand{}

\newcommand{\EndExtendedAbstract}
{
  \EndOfPage

  % 設定段落之間的距離
  \setlength{\parskip}{0.3cm}

  % Reset page
  \setlength{\parindent}{1.5em}
  \baselineskip=26pt

  % Re-set style
  \pagestyle{plain}
} % End of \def{}

% Summary in Extended Abstract
\newcommand{\ExtAbstractSummary}[1]{
  \begin{framed}
    \ExtAbstractChapter{SUMMARY}{#1}

    % Keyword
    \ExtendedAbstractKeyword
  \end{framed}
} % End of \newcommand{}

% Chapter in Extended Abstract
\newcommand{\ExtAbstractChapter}[2]
{
  \vspace{0.4cm}
  \centerline{\textbf{#1}}

  #2
} % End of \newcommand{}

% Sub-chapter in Extended Abstract
\newcommand{\ExtAbstractSection}[2]
{
  \textbf{#1}

  #2
} % End of \newcommand{}

\newcommand{\AbstractKeyword}[0]
{
  % Keyword
  \par
  \ifthenelse{\equal{\GetKeywordsE}{\empty}}
  {
    \ifthenelse{\equal{\GetKeywordsD}{\empty}}
    {
      \ifthenelse{\equal{\GetKeywordsC}{\empty}}
      {
        \ifthenelse{\equal{\GetKeywordsB}{\empty}}
        {
          \ifthenelse{\equal{\GetKeywordsA}{\empty}}
          {}
          {
            {\noindent \bf 關鍵字:} \GetKeywordsA
          } % End of else{}
        } % End of if{}
        {
          {\noindent \bf 關鍵字:} \GetKeywordsA, \GetKeywordsB
        } % End of else{}
      } % End of if{}
      {
        {\noindent \bf 關鍵字:} \GetKeywordsA, \GetKeywordsB, \GetKeywordsC
      } % End of else{}
    } % End of if{}
    {
      {\noindent \bf 關鍵字:} \GetKeywordsA, \GetKeywordsB, \GetKeywordsC, \GetKeywordsD
    } % End of else{}
  } % End of if{}
  {
    {\noindent\bf 關鍵字:} \GetKeywordsA, \GetKeywordsB, \GetKeywordsC, \GetKeywordsD, \GetKeywordsE
  } % End of else{}
} % End of \newcommand{}

\newcommand{\ExtendedAbstractKeyword}[0]
{
  % Keyword
  \par
  \ifthenelse{\equal{\GetKeywordsE}{\empty}}
  {
    \ifthenelse{\equal{\GetKeywordsD}{\empty}}
    {
      \ifthenelse{\equal{\GetKeywordsC}{\empty}}
      {
        \ifthenelse{\equal{\GetKeywordsB}{\empty}}
        {
          \ifthenelse{\equal{\GetKeywordsA}{\empty}}
          {}
          {
            {\noindent \bf Key words:} \GetKeywordsA
          } % End of else{}
        } % End of if{}
        {
          {\noindent \bf Key words:} \GetKeywordsA, \GetKeywordsB
        } % End of else{}
      } % End of if{}
      {
        {\noindent \bf Key words:} \GetKeywordsA, \GetKeywordsB, \GetKeywordsC
      } % End of else{}
    } % End of if{}
    {
      {\noindent \bf Key words:} \GetKeywordsA, \GetKeywordsB, \GetKeywordsC, \GetKeywordsD
    } % End of else{}
  } % End of if{}
  {
    {\noindent\bf Key words:} \GetKeywordsA, \GetKeywordsB, \GetKeywordsC, \GetKeywordsD, \GetKeywordsE
  } % End of else{}
} % End of \newcommand{}

% ----------------------------------------------------------------------------


%
% This file is part of ncku-thesis-templete.
%
% ncku-thesis-templete is distributed in the hope that it will be useful,
% you can redistribute it and/or modify
% it under the terms of the Attribution-NonCommercial-ShareAlike
% 4.0 International.
%
% You should have received a copy of the
% Attribution-NonCommercial-ShareAlike 4.0 International
% along with ncku-thesis-templete.
%
% If not, see <http://creativecommons.org/licenses/by-nc-sa/4.0/legalcode.txt>.
%

% Some helper function about acknowledgments

% ----------------------------------------------------------------------------

% Acknowledgments
\newcommand{\StartAcknowledgments}
{
  \StartNewPage

  % Add to "Table of Contents"
  \addcontentsline{toc}{chapter}{Acknowledgments}

  % Title
  \centerline{\Large \textbf{Acknowledgments}}

%  \IfNoValueF{#1}{\label{#1}}
} % End of \newcommand{}

%\newcommand{\StartAcknowledgmentsChi}[1]
\newcommand{\StartAcknowledgmentsChi}
{
  \StartNewPage

  % Add to "Table of Contents"
  \addcontentsline{toc}{chapter}{致謝}

  % Title
  \centerline{\Large \textbf{致謝}}

%  \IfNoValueF{#1}{\label{#1}}
} % End of \newcommand{}

\def \EndAcknowledgments {\EndOfPage}

% ----------------------------------------------------------------------------


%
% This file is part of ncku-thesis-templete.
%
% ncku-thesis-templete is distributed in the hope that it will be useful,
% you can redistribute it and/or modify
% it under the terms of the Attribution-NonCommercial-ShareAlike
% 4.0 International.
%
% You should have received a copy of the
% Attribution-NonCommercial-ShareAlike 4.0 International
% along with ncku-thesis-templete.
%
% If not, see <http://creativecommons.org/licenses/by-nc-sa/4.0/legalcode.txt>.
%

% Some helper function about appendix

% ----------------------------------------------------------------------------

\def \StartAppendix
{
  \StartNewPage
  % Appendix begins here (For one appendix chapter only)
  %\cleardoublepage
  \appendix
  %\clearpage
  %\setcounter{page}{\thepage}
  % If more than one appendix chapters
  %\appendices
}

\def \EndAppendix {\EndOfPage}

% ----------------------------------------------------------------------------


%
% This file is part of ncku-thesis-templete.
%
% ncku-thesis-templete is distributed in the hope that it will be useful,
% you can redistribute it and/or modify
% it under the terms of the Attribution-NonCommercial-ShareAlike
% 4.0 International.
%
% You should have received a copy of the
% Attribution-NonCommercial-ShareAlike 4.0 International
% along with ncku-thesis-templete.
%
% If not, see <http://creativecommons.org/licenses/by-nc-sa/4.0/legalcode.txt>.
%

% Some helper function about references and bibliography
% 參考文獻或資料

% ----------------------------------------------------------------------------

\DeclareDocumentCommand{\ReferencesFiles}{%
  G{\empty} G{\empty} G{\empty} %
  G{\empty} G{\empty} G{\empty} %
  G{\empty} G{\empty} G{\empty}}
{
  \StartNewPage

  % Add to "Table of Contents"
  %\addcontentsline{toc}{chapter}{References}

  % Change the title of bibliography
  %\renewcommand\bibname{\centerline{\Large References}}

  % ------------------------------------------------

  %Style Name | Author Name Format | Reference Format | Sorting
  %  abbrv    | H. J. Simpson      |        #ID#      | by author
  \bibliographystyle{abbrv}

  % ------------------------------------------------

  % Bib files
  \ifthenelse{\equal{#9}{\empty}}
  {
    \ifthenelse{\equal{#8}{\empty}}
    {
      \ifthenelse{\equal{#7}{\empty}}
      {
        \ifthenelse{\equal{#6}{\empty}}
        {
          \ifthenelse{\equal{#5}{\empty}}
          {
            \ifthenelse{\equal{#4}{\empty}}
            {
              \ifthenelse{\equal{#3}{\empty}}
              {
                \ifthenelse{\equal{#2}{\empty}}
                {
                  \ifthenelse{\equal{#1}{\empty}}
                  {} % End of if{}
                  {\bibliography{#1}} % End of else{}
                } % End of if{}
                {
                  \bibliography{#1,#2}
                } % End of else{}
              } % End of if{}
              {
                \bibliography{#1,#2,#3}
              } % End of else{}
            } % End of if{}
            {
              \bibliography{#1,#2,#3,#4}
            } % End of else{}
          } % End of if{}
          {
            \bibliography{#1,#2,#3,#4,#5}
          } % End of else{}
        } % End of if{}
        {
          \bibliography{#1,#2,#3,#4,#5,#6}
        } % End of else{}
      } % End of if{}
      {
        \bibliography{#1,#2,#3,#4,#5,#6,#7}
      } % End of else{}
    } % End of if{}
    {
      \bibliography{#1,#2,#3,#4,#5,#6,#7,#8}
    } % End of else{}
  } % End of if{}
  {
    \bibliography{#1,#2,#3,#4,#5,#6,#7,#8,#9}
  } % End of else{}

  % ------------------------------------------------

  \EndOfPage

} % End of \DeclareDocumentCommand{}
% ----------------------------------------------------------------------------


%
% This file is part of ncku-thesis-templete.
%
% ncku-thesis-templete is distributed in the hope that it will be useful,
% you can redistribute it and/or modify
% it under the terms of the Attribution-NonCommercial-ShareAlike
% 4.0 International.
%
% You should have received a copy of the
% Attribution-NonCommercial-ShareAlike 4.0 International
% along with ncku-thesis-templete.
%
% If not, see <http://creativecommons.org/licenses/by-nc-sa/4.0/legalcode.txt>.
%

% Some helper function use in index

% ----------------------------------------------------------------------------

% Use to include the index files
\newcommand{\DisplayIndex}{%
% This file is part of ncku-thesis-templete.
%
% ncku-thesis-templete is distributed in the hope that it will be useful,
% you can redistribute it and/or modify
% it under the terms of the Attribution-NonCommercial-ShareAlike
% 4.0 International.
%
% You should have received a copy of the
% Attribution-NonCommercial-ShareAlike 4.0 International
% along with ncku-thesis-templete.
%
% If not, see <http://creativecommons.org/licenses/by-nc-sa/4.0/legalcode.txt>.
%

% 目錄 (內容, 圖表和圖片)
% Index of contents, tables and figures
\newpage

\phantomsection
\addcontentsline{toc}{chapter}{Table of Contents}
\tableofcontents

\phantomsection
%\addcontentsline{toc}{chapter}{Tables}
\listoftables

\phantomsection
%\addcontentsline{toc}{chapter}{Figures} 
\listoffigures

}

% ----------------------------------------------------------------------------


% Some function that use for school
%
% This file is part of ncku-thesis-templete.
%
% ncku-thesis-templete is distributed in the hope that it will be useful,
% you can redistribute it and/or modify
% it under the terms of the Attribution-NonCommercial-ShareAlike
% 4.0 International.
%
% You should have received a copy of the
% Attribution-NonCommercial-ShareAlike 4.0 International
% along with ncku-thesis-templete.
%
% If not, see <http://creativecommons.org/licenses/by-nc-sa/4.0/legalcode.txt>.
%

% ----------------------------------------------------------------------------
% The list of all college in NCKU

% Use the list from:
%   http://web.ncku.edu.tw/files/11-1000-182.php

% 所有學院跟系所的設定
% 縮寫是靠學校給的Domain name所得出的, 故可能會有錯誤的時候
% ----------------------------------------------------------------------------

% ----------------------------------------------------------------------------
% 學院 College
% ----------------------------------------------------------------------------

% 文學院 College of Liberal Arts
\newcommand{\CollLiberalArts}[0]
{
  \SetCollName{文學院}{College of Liberal Arts}
} % End of \newcommand{}

% 理學院 College of Sciences
\newcommand{\CollSciences}[0]
{
  \SetCollName{理學院}{College of Sciences}
} % End of \newcommand{}

% 工學院 College of Engineering
\newcommand{\CollEngineering}[0]
{
  \SetCollName{工學院}{College of Engineering}
} % End of \newcommand{}

% 電機資訊學院 College of Electrical Engineering & Computer Science
\newcommand{\CollElectricalEngineeringAndComputerScience}[0]
{
  \SetCollName{電機資訊學院}{College of Electrical Engineering and Computer Science}
} % End of \newcommand{}

% 管理學院 College of Management
\newcommand{\CollManagement}[0]
{
  \SetCollName{管理學院}{College of Management}
} % End of \newcommand{}

% 社會科學院 College of Social Science
\newcommand{\CollSocialScience}[0]
{
  \SetCollName{社會科學院}{College of Social Science}
} % End of \newcommand{}

% 規劃與設計學院 College of Planning & Design
\newcommand{\CollPlanningAndDesign}[0]
{
  \SetCollName{規劃與設計學院}{College of Planning and Design}
} % End of \newcommand{}

% 生物科學與科技學院 College of Bioscience & Biotechnology
\newcommand{\CollBioscienceAndBiotechnology}[0]
{
  \SetCollName{生物科學與科技學院}{College of Bioscience and Biotechnology}
} % End of \newcommand{}

% 醫學院 College of Medicine
\newcommand{\CollMedicine}[0]
{
  \SetCollName{醫學院}{College of Medicine}
} % End of \newcommand{}

% ----------------------------------------------------------------------------


%
% This file is part of ncku-thesis-templete.
%
% ncku-thesis-templete is distributed in the hope that it will be useful,
% you can redistribute it and/or modify
% it under the terms of the Attribution-NonCommercial-ShareAlike
% 4.0 International.
%
% You should have received a copy of the
% Attribution-NonCommercial-ShareAlike 4.0 International
% along with ncku-thesis-templete.
%
% If not, see <http://creativecommons.org/licenses/by-nc-sa/4.0/legalcode.txt>.
%

% ----------------------------------------------------------------------------
% The list of all department and institute in NCKU

% Use the list from:
%   http://web.ncku.edu.tw/files/11-1000-182.php

% 所有學院跟系所的設定.
%
% 縮寫是靠學校給的Domain name所得出的, 故可能會有錯誤的時候.
% 所以如果錯了的話, 就請告知真正的寫法(或縮寫)是什麼.
%
% ----------------------------------------------------------------------------


% ----------------------------------------------------------------------------
%                       系所 Department and Institute 
% ----------------------------------------------------------------------------


% ----------------------------------------------------------------------------
% --------------------- 文學院 College of Liberal Arts ---------------------
% ----------------------------------------------------------------------------

% 中國文學系 Department of Chinese Literature
\newcommand{\SetDeptChinese}[0]
{
  \SetDeptName{中國文學系}{Chinese}{Department of Chinese Literature}
  \CollLiberalArts
} % End of \newcommand{}

% 藝術研究所 Institute of Art
\newcommand{\SetDeptArt}[0]
{
  \SetDeptName{藝術研究所}{Art}{Institute of Art}
  \CollLiberalArts
} % End of \newcommand{}

% 閩南文化研究中心 Min-Nan Culture Studies Center
\newcommand{\SetDeptMinNan}[0]
{
  \SetDeptName{閩南文化研究中心}{MinNan}{Min-Nan Culture Studies Center}
  \CollLiberalArts
} % End of \newcommand{}

% 外國語文學系 Department of Foreign Languages and Literature
\newcommand{\SetDeptFLLD}[0]
{
  \SetDeptName{外國語文學系}{FLLD}{Department of Foreign Languages and Literature}
  \CollLiberalArts
} % End of \newcommand{}

% 臺灣文學系 Department of Taiwanese Literature
\newcommand{\SetDeptTWL}[0]
{
  \SetDeptName{臺灣文學系}{TWL}{Department of Taiwanese Literature}
  \CollLiberalArts
} % End of \newcommand{}

% 華語中心 Chinese Language Center
\newcommand{\SetDeptKCLC}[0]
{
  \SetDeptName{華語中心}{KCLC}{Chinese Language Center}
  \CollLiberalArts
} % End of \newcommand{}

% 外語中心 Foreign Language Center
\newcommand{\SetDeptLang}[0]
{
  \SetDeptName{外語中心}{Lang}{Foreign Language Center}
  \CollLiberalArts
} % End of \newcommand{}

% 歷史學系 Department of History
\newcommand{\SetDeptHis}[0]
{
  \SetDeptName{歷史學系}{His}{Department of History}
  \CollLiberalArts
} % End of \newcommand{}

% ----------------------------------------------------------------------------
% --------------------- 理學院 College of Sciences ---------------------
% ----------------------------------------------------------------------------

% 數學系 Department of Mathematics
\newcommand{\SetDeptMath}[0]
{
  \SetDeptName{數學系}{Math}{Department of Mathematics}
  \CollSciences
} % End of \newcommand{}

% 光電科學與工程學系 Departmment of Photonics
\newcommand{\SetDeptDPS}[0]
{
  \SetDeptName{光電科學與工程學系}{DPS}{Departmment of Photonics}
  \CollSciences
} % End of \newcommand{}

% 物理學系 Department of Physics
\newcommand{\SetDeptPhys}[0]
{
  \SetDeptName{物理學系}{Phys}{Department of Physics}
  \CollSciences
} % End of \newcommand{}

% 化學系 Department of Chemistry
\newcommand{\SetDeptCh}[0]
{
  \SetDeptName{化學系}{Ch}{Department of Chemistry}
  \CollSciences
} % End of \newcommand{}

% 地球科學系 Department of Earth Sciences
\newcommand{\SetDeptEarth}[0]
{
  \SetDeptName{地球科學系}{Earth}{Department of Earth Sciences}
  \CollSciences
} % End of \newcommand{}

% 太空與電漿科學研究所 Institute of Space and Plasma Sciences
\newcommand{\SetDeptPSSC}[0]
{
  \SetDeptName{太空與電漿科學研究所}{PSSC}{Institute of Space and Plasma Sciences}
  \CollSciences
} % End of \newcommand{}

% 國家理論科學研究中心 National Center for Theoretical Sciences (South)
\newcommand{\SetDeptNCTS}[0]
{
  \SetDeptName{國家理論科學研究中心}{NCTS}{National Center for Theoretical Sciences (South)}
  \CollSciences
} % End of \newcommand{}

% ----------------------------------------------------------------------------
% --------------------- 工學院 College of Engineering ---------------------
% ----------------------------------------------------------------------------

% 機械工程學系 Department of Mechanical Engineering
\newcommand{\SetDeptME}[0]
{
  \SetDeptName{機械工程學系}{ME}{Department of Mechanical Engineering}
  \CollEngineering
} % End of \newcommand{}

% 化學工程學系 Department of Chemical Engineering
\newcommand{\SetDeptChe}[0]
{
  \SetDeptName{化學工程學系}{Che}{Department of Chemical Engineering}
  \CollEngineering
} % End of \newcommand{}

% 土木工程學系 Department of Civil Engineering
\newcommand{\SetDeptCivil}[0]
{
  \SetDeptName{土木工程學系}{Civil}{Department of Civil Engineering}
  \CollEngineering
} % End of \newcommand{}

% 材料科學及工程學系 Department of Materials Science and Engineering
\newcommand{\SetDeptMSE}[0]
{
  \SetDeptName{材料科學及工程學系}{MSE}{Department of Materials Science and Engineering}
  \CollEngineering
} % End of \newcommand{}

% 水利及海洋工程學系 Department of Hydraulic and Ocean Engineering
\newcommand{\SetDeptHyd}[0]
{
  \SetDeptName{水利及海洋工程學系}{Hyd}{Department of Hydraulic and Ocean Engineering}
  \CollEngineering
} % End of \newcommand{}

% 工程科學系 Department of Engineering Science
\newcommand{\SetDeptES}[0]
{
  \SetDeptName{工程科學系}{ES}{Department of Engineering Science}
  \CollEngineering
} % End of \newcommand{}

% 系統及船舶機電工程學系 Department of System and Naval Mechatronic Engineering
\newcommand{\SetDeptSNAME}[0]
{
  \SetDeptName{系統及船舶機電工程學系}{SNAME}{Department of System and Naval Mechatronic Engineering}
  \CollEngineering
} % End of \newcommand{}

% 航空太空工程學系 Department of Aeronautics and Astronautics
\newcommand{\SetDeptIAA}[0]
{
  \SetDeptName{航空太空工程學系}{IAA}{Department of Aeronautics and Astronautics}
  \CollEngineering
} % End of \newcommand{}

% 資源工程學系 Department of Resources Engineering
\newcommand{\SetDeptMP}[0]
{
  \SetDeptName{資源工程學系}{MP}{Department of Resources Engineering}
  \CollEngineering
} % End of \newcommand{}

% 環境工程學系 Department of Environmental Engineering
\newcommand{\SetDeptEV}[0]
{
  \SetDeptName{環境工程學系}{EV}{Department of Environmental Engineering}
  \CollEngineering
} % End of \newcommand{}

% 生物醫學工程學系 Department of BioMedical Engineering
\newcommand{\SetDeptBME}[0]
{
  \SetDeptName{生物醫學工程學系}{BME}{Department of BioMedical Engineering}
  \CollEngineering
} % End of \newcommand{}

% 測量及空間資訊學系 Department of Geomatics
\newcommand{\SetDeptGeomatics}[0]
{
  \SetDeptName{測量及空間資訊學系}{Geomatics}{Department of Geomatics}
  \CollEngineering
} % End of \newcommand{}

% 海洋科技與事務研究所 Institute of Ocean Technology and Marine Affairs
\newcommand{\SetDeptIOTMA}[0]
{
  \SetDeptName{海洋科技與事務研究所}{IOTMA}{Institute of Ocean Technology and Marine Affairs}
  \CollEngineering
} % End of \newcommand{}

% 民航研究所 Institute of Civil Aviation
\newcommand{\SetDeptICA}[0]
{
  \SetDeptName{民航研究所}{ICA}{Institute of Civil Aviation}
  \CollEngineering
} % End of \newcommand{}

% 能源國際學士學位學程 International Bachelor Degree Program on Energy
\newcommand{\SetDeptIBDPE}[0]
{
  \SetDeptName{能源國際學士學位學程}{IBDPE}{International Bachelor Degree Program on Energy}
  \CollEngineering
} % End of \newcommand{}

% 尖端材料國際碩士學位學程 International Curriculum for Advanced Materials Program
\newcommand{\SetDeptICAMP}[0]
{
  \SetDeptName{尖端材料國際碩士學位學程}{ICAMP}{International Curriculum for Advanced Materials Program}
  \CollEngineering
} % End of \newcommand{}

% 自然災害減災及管理國際碩士學位學程 International Master Program on Natural Hazards Mitigation and Management
\newcommand{\SetDeptINHMM}[0]
{
  \SetDeptName{自然災害減災及管理國際碩士學位學程}{INHMM}{International Master Program on \\ Natural Hazards Mitigation and Management}
  \CollEngineering
} % End of \newcommand{}

% 工程管理碩士在職專班 International Graduate Program of Civil Engineering and Management
\newcommand{\SetDeptICEM}[0]
{
  \SetDeptName{工程管理碩士在職專班}{ICEM}{International Graduate Program of \\ Civil Engineering and Management}
  \CollEngineering
} % End of \newcommand{}

% ----------------------------------------------------------------------------
% ------ 電機資訊學院 College of Electrical Engineering & Computer Science ------
% ----------------------------------------------------------------------------

% 電機工程學系 Department of Electrical Engineering
\newcommand{\SetDeptEE}[0]
{
  \SetDeptName{電機工程學系}{EE}{Department of Electrical Engineering}
  \CollElectricalEngineeringAndComputerScience
} % End of \newcommand{}

% 資訊工程學系 Insitute of Computer Science and Information Engineering
\newcommand{\SetDeptCSIE}[0]
{
  \SetDeptName{資訊工程學系}{CSIE}{Insitute of Computer Science and Information Engineering}
  \CollElectricalEngineeringAndComputerScience
} % End of \newcommand{}

% 微電子工程研究所 Institute of Microelectronics
\newcommand{\SetDeptIME}[0]
{
  \SetDeptName{微電子工程研究所}{IME}{Institute of Microelectronics}
  \CollElectricalEngineeringAndComputerScience
} % End of \newcommand{}

% 電腦與通信工程研究所 Institute of Computer & Communication Engineering
\newcommand{\SetDeptCCE}[0]
{
  \SetDeptName{電腦與通信工程研究所}{CCE}{Institute of Computer & Communication Engineering}
  \CollElectricalEngineeringAndComputerScience
} % End of \newcommand{}

% 製造資訊與系統研究所 Institute of Manufacturing Information and Systems
\newcommand{\SetDeptIMIS}[0]
{
  \SetDeptName{製造資訊與系統研究所}{IMIS}{Institute of Manufacturing Information and Systems}
  \CollElectricalEngineeringAndComputerScience
} % End of \newcommand{}

% 醫學資訊研究所 Institute of Medical Informatics
\newcommand{\SetDeptIMI}[0]
{
  \SetDeptName{醫學資訊研究所}{IMI}{Institute of Medical Informatics}
  \CollElectricalEngineeringAndComputerScience
} % End of \newcommand{}

% ----------------------------------------------------------------------------
% --------------------- 管理學院 College of Management ---------------------
% ----------------------------------------------------------------------------

% 統計學系 Department of Statistics
\newcommand{\SetDeptSTAT}[0]
{
  \SetDeptName{統計學系}{STAT}{Department of Statistics}
  \CollManagement
} % End of \newcommand{}

% 會計學系 Department of Accountancy
\newcommand{\SetDeptACC}[0]
{
  \SetDeptName{會計學系}{ACC}{Department of Accountancy}
  \CollManagement
} % End of \newcommand{}

% 交通管理科學系 Department of Transportation and Communication Management Science
\newcommand{\SetDeptTCM}[0]
{
  \SetDeptName{交通管理科學系}{TCM}{Department of Transportation and Communication Management Science}
  \CollManagement
} % End of \newcommand{}

% 企業管理學系暨國際企業研究所 Department of Business Administration and Graduate Institute of International Business
\newcommand{\SetDeptBA}[0]
{
  \SetDeptName{企業管理學系暨國際企業研究所}{BA}{Department of Business Administration and \\ Graduate Institute of International Business}
  \CollManagement
} % End of \newcommand{}

% 電信管理研究所 Institute of Telecommunications Management
\newcommand{\SetDeptTM}[0]
{
  \SetDeptName{電信管理研究所}{TM}{Institute of Telecommunications Management}
  \CollManagement
} % End of \newcommand{}

% 工業與資訊管理學系暨資訊管理研究所 Institute of Information Management
\newcommand{\SetDeptIIM}[0]
{
  \SetDeptName{工業與資訊管理學系暨資訊管理研究所}{IIM}{Institute of Information Management}
  \CollManagement
} % End of \newcommand{}

% 財務金融研究所 Institute of Finance & Banking
\newcommand{\SetDeptFin}[0]
{
  \SetDeptName{財務金融研究所}{Fin}{Institute of Finance & Banking}
  \CollManagement
} % End of \newcommand{}

% 體育健康與休閒研究所 Institute of Physical Education, Health & Leisure Studies
\newcommand{\SetDeptPHEI}[0]
{
  \SetDeptName{體育健康與休閒研究所}{PHEI}{Institute of Physical Education, \\ Health & Leisure Studies}
  \CollManagement
} % End of \newcommand{}

% 高階管理碩士在職專班 (EMBA) Executive Master of Business Administration
\newcommand{\SetDeptEMBA}[0]
{
  \SetDeptName{高階管理碩士在職專班}{EMBA}{Executive Master of Business Administration}
  \CollManagement
} % End of \newcommand{}

% 國際經營管理研究所 (IMBA) Institute of International Management (IMBA)
\newcommand{\SetDeptIMBA}[0]
{
  \SetDeptName{國際經營管理研究所}{IMBA}{Institute of International Management}
  \CollManagement
} % End of \newcommand{}

% 經營管理碩士班 (AMBA) Advanced Master of Business Administration
\newcommand{\SetDeptAMBA}[0]
{
  \SetDeptName{經營管理碩士班}{AMBA}{Advanced Master of Business Administration}
  \CollManagement
} % End of \newcommand{}

% ----------------------------------------------------------------------------
% ------------------- 社會科學院 College of Social Science -------------------
% ----------------------------------------------------------------------------

% 政治學系 Department of Political Science
\newcommand{\SetDeptPolSci}[0]
{
  \SetDeptName{政治學系}{PolSci}{Department of Political Science}
  \CollSocialScience
} % End of \newcommand{}

% 經濟學系 Department of Economics
\newcommand{\SetDeptEconomic}[0]
{
  \SetDeptName{經濟學系}{Economic}{Department of Economics}
  \CollSocialScience
} % End of \newcommand{}

% 心理學系 Department of Psychology
\newcommand{\SetDeptPsychology}[0]
{
  \SetDeptName{心理學系}{Psychology}{Department of Psychology}
  \CollSocialScience
} % End of \newcommand{}

% 法律學系 Department of Law and Institute of Law in Science and Technology
\newcommand{\SetDeptLaw}[0]
{
  \SetDeptName{法律學系}{Law}{Department of Law and \\Institute of Law in Science and Technology}
  \CollSocialScience
} % End of \newcommand{}

% 教育研究所 Institute of Education
\newcommand{\SetDeptED}[0]
{
  \SetDeptName{教育研究所}{ED}{Institute of Education}
  \CollSocialScience
} % End of \newcommand{}

% 認知科學研究所 Institute of Cognitive Science
\newcommand{\SetDeptIOCS}[0]
{
  \SetDeptName{認知科學研究所}{IOCS}{Institute of Cognitive Science}
  \CollSocialScience
} % End of \newcommand{}

% 政治經濟學研究所 Institute of Political Economy
\newcommand{\SetDeptGIPE}[0]
{
  \SetDeptName{政治經濟學研究所}{GIPE}{Institute of Political Economy}
  \CollSocialScience
} % End of \newcommand{}

% 心智影像研究中心 Mind Research and Image Center
\newcommand{\SetDeptFMRI}[0]
{
  \SetDeptName{心智影像研究中心}{FMRI}{Mind Research and Image Center}
  \CollSocialScience
} % End of \newcommand{}

% ----------------------------------------------------------------------------
% ---------------- 規劃與設計學院 College of Planning & Design ----------------
% ----------------------------------------------------------------------------

% 建築學系 Department of Architecture
\newcommand{\SetDeptArch}[0]
{
  \SetDeptName{建築學系}{Arch}{Department of Architecture}
  \CollPlanningAndDesign
} % End of \newcommand{}

% 都市計劃學系 Department of Urban Planning
\newcommand{\SetDeptUP}[0]
{
  \SetDeptName{都市計劃學系}{UP}{Department of Urban Planning}
  \CollPlanningAndDesign
} % End of \newcommand{}

% 工業設計學系 Department of Industrial Design
\newcommand{\SetDeptID}[0]
{
  \SetDeptName{工業設計學系}{ID}{Department of Industrial Design}
  \CollPlanningAndDesign
} % End of \newcommand{}

% 創意產業設計研究所 Institute of Creative Industry Design
\newcommand{\SetDeptICID}[0]
{
  \SetDeptName{創意產業設計研究所}{ICID}{Institute of Creative Industry Design}
  \CollPlanningAndDesign
} % End of \newcommand{}

% ----------------------------------------------------------------------------
% ---------- 生物科學與科技學院 College of Bioscience & Biotechnology ----------
% ----------------------------------------------------------------------------

% 生命科學系 Department of Life Sciences
\newcommand{\SetDeptBio}[0]
{
  \SetDeptName{生命科學系}{Bio}{Department of Life Sciences}
  \CollBioscienceAndBiotechnology
} % End of \newcommand{}

% 生物科技研究所 Institute of Biotechnology
\newcommand{\SetDeptBioTech}[0]
{
  \SetDeptName{生物科技研究所}{BioTech}{Institute of Biotechnology}
  \CollBioscienceAndBiotechnology
} % End of \newcommand{}

% 生物資訊與訊息傳遞研究所 Institute of Bioinformatics and Biosignal Transduction
\newcommand{\SetDeptIBBT}[0]
{
  \SetDeptName{生物資訊與訊息傳遞研究所}{IBBT}{Institute of Bioinformatics and \\ Biosignal Transduction}
  \CollBioscienceAndBiotechnology
} % End of \newcommand{}

% 熱帶植物科學研究所 Institute of Tropical Plant Sciences
\newcommand{\SetDeptITPS}[0]
{
  \SetDeptName{熱帶植物科學研究所}{ITPS}{Institute of Tropical Plant Sciences}
  \CollBioscienceAndBiotechnology
} % End of \newcommand{}

% ----------------------------------------------------------------------------
% --------------------- 醫學院 College of Medicine ---------------------
% ----------------------------------------------------------------------------

% 醫學系 School of Medicine
\newcommand{\SetDeptEDUC}[0]
{
  \SetDeptName{醫學系}{EDUC}{School of Medicine}
  \CollMedicine
} % End of \newcommand{}

% 生物化學暨分子生物學研究所 Department of Biochemistry and Molecular Biology
\newcommand{\SetDeptBiohem}[0]
{
  \SetDeptName{生物化學暨分子生物學研究所}{Biohem}{Department of Biochemistry and \\ Molecular Biology}
  \CollMedicine
} % End of \newcommand{}

% 病理學科 Department of Pathology
\newcommand{\SetDeptPath}[0]
{
  \SetDeptName{病理學科}{Path}{Department of Pathology}
  \CollMedicine
} % End of \newcommand{}

% 內科學科 Department of Internal Medicine
\newcommand{\SetDeptIntMed}[0]
{
  \SetDeptName{內科學科}{IntMed}{Department of Internal Medicine}
  \CollMedicine
} % End of \newcommand{}

% 生理學研究所 Department of Physiology
\newcommand{\SetDeptPhysMed}[0]
{
  \SetDeptName{生理學研究所}{PhysMed}{Department of Physiology}
  \CollMedicine
} % End of \newcommand{}

% 外科學科 Department of Surgery
\newcommand{\SetDeptSurgery}[0]
{
  \SetDeptName{外科學科}{Surgery}{Department of Surgery}
  \CollMedicine
} % End of \newcommand{}

% 小兒學科 Department of Pediatrics
\newcommand{\SetDeptPed}[0]
{
  \SetDeptName{小兒學科}{Ped}{Department of Pediatrics}
  \CollMedicine
} % End of \newcommand{}

% 解剖學科暨細胞生物與解剖學研究所 Department of Cell Biology and Anatomy
\newcommand{\SetDeptAnatomy}[0]
{
  \SetDeptName{解剖學科暨細胞生物與解剖學研究所}{Anatomy}{Department of Cell Biology and Anatomy}
  \CollMedicine
} % End of \newcommand{}

% 婦產學科 Department of Obstetrics and Gynecology
\newcommand{\SetDeptObsGyn}[0]
{
  \SetDeptName{婦產學科}{ObsGyn}{Department of Obstetrics and Gynecology}
  \CollMedicine
} % End of \newcommand{}

% 骨科學科 Department of Orthopaedics
\newcommand{\SetDeptBone}[0]
{
  \SetDeptName{骨科學科}{Bone}{Department of Orthopaedics}
  \CollMedicine
} % End of \newcommand{}

% 公共衛生學科暨公共衛生研究所 Department of Public Health
\newcommand{\SetDeptPhMed}[0]
{
  \SetDeptName{公共衛生學科暨公共衛生研究所}{PhMed}{Department of Public Health}
  \CollMedicine
} % End of \newcommand{}

% 神經學科 Department of Neurology
\newcommand{\SetDeptNeuro}[0]
{
  \SetDeptName{神經學科}{Neuro}{Department of Neurology}
  \CollMedicine
} % End of \newcommand{}

% 精神學科 Department of Psychiatry
\newcommand{\SetDeptPsy}[0]
{
  \SetDeptName{精神學科}{Psy}{Department of Psychiatry}
  \CollMedicine
} % End of \newcommand{}

% 寄生蟲學科 Department of Parasitology
\newcommand{\SetDeptParasite}[0]
{
  \SetDeptName{寄生蟲學科}{Parasite}{Department of Parasitology}
  \CollMedicine
} % End of \newcommand{}

% 眼科學科 Department of Ophthalmology
\newcommand{\SetDeptOphth}[0]
{
  \SetDeptName{眼科學科}{Ophth}{Department of Ophthalmology}
  \CollMedicine
} % End of \newcommand{}

% 耳鼻喉學科 Department of Otolaryngology
\newcommand{\SetDeptOtolaryngo}[0]
{
  \SetDeptName{耳鼻喉學科}{Otolaryngo}{Department of Otolaryngology}
  \CollMedicine
} % End of \newcommand{}

% 工業衛生學科暨環境醫學研究所 Department of Environmental and Occupational Health
\newcommand{\SetDeptDEOH}[0]
{
  \SetDeptName{工業衛生學科暨環境醫學研究所}{DEOH}{Department of Environmental and Occupational Health}
  \CollMedicine
} % End of \newcommand{}

% 皮膚學科 Department of Dermatology
\newcommand{\SetDeptDerm}[0]
{
  \SetDeptName{皮膚學科}{Derm}{Department of Dermatology}
  \CollMedicine
} % End of \newcommand{}

% 泌尿學科 Department of Urology
\newcommand{\SetDeptUro}[0]
{
  \SetDeptName{泌尿學科}{Uro}{Department of Urology}
  \CollMedicine
} % End of \newcommand{}

% 藥理學科暨藥理學研究所 Department of Pharmacology
\newcommand{\SetDeptPharmaco}[0]
{
  \SetDeptName{藥理學科暨藥理學研究所}{Pharmaco}{Department of Pharmacology}
  \CollMedicine
} % End of \newcommand{}

% 麻醉學科 Department of Anesthesiology
\newcommand{\SetDeptAnesth}[0]
{
  \SetDeptName{麻醉學科}{Anesth}{Department of Anesthesiology}
  \CollMedicine
} % End of \newcommand{}

% 復健學科 Department of Physical Medicine and Rehabilitation
\newcommand{\SetDeptRehab}[0]
{
  \SetDeptName{復健學科}{Rehab}{Department of Physical Medicine and Rehabilitation}
  \CollMedicine
} % End of \newcommand{}

% 微生物學及免疫研究所 Department of Microbiology and Immunology
\newcommand{\SetDeptMicrobio}[0]
{
  \SetDeptName{微生物學及免疫研究所}{Microbio}{Department of Microbiology and Immunology}
  \CollMedicine
} % End of \newcommand{}

% 放射線學科 Department of Diagnostic Radiology
\newcommand{\SetDeptRad}[0]
{
  \SetDeptName{放射線學科}{Rad}{Department of Diagnostic Radiology}
  \CollMedicine
} % End of \newcommand{}

% 核子醫學科 Department of Nuclear Medicine
\newcommand{\SetDeptNM}[0]
{
  \SetDeptName{核子醫學科}{NM}{Department of Nuclear Medicine}
  \CollMedicine
} % End of \newcommand{}

% 家庭醫學科 Department of Family Medicine
\newcommand{\SetDeptFamily}[0]
{
  \SetDeptName{家庭醫學科}{Family}{Department of Family Medicine}
  \CollMedicine
} % End of \newcommand{}

% 急診學科 Department of Emergency Medicine
\newcommand{\SetDeptEmergency}[0]
{
  \SetDeptName{急診學科}{Emergency}{Department of Emergency Medicine}
  \CollMedicine
} % End of \newcommand{}

% 牙科學科 Department of Dentistry
\newcommand{\SetDeptDentistry}[0]
{
  \SetDeptName{牙科學科}{Dentistry}{Department of Dentistry}
  \CollMedicine
} % End of \newcommand{}

% 職業及環境醫學科 Department of Occupational and Environmental Medicine
\newcommand{\SetDeptOEM}[0]
{
  \SetDeptName{職業及環境醫學科}{OEM}{Department of Occupational and Environmental Medicine}
  \CollMedicine
} % End of \newcommand{}

% 法醫學科 Department of Forensic Medicine
\newcommand{\SetDeptForensic}[0]
{
  \SetDeptName{法醫學科}{Forensic}{Department of Forensic Medicine}
  \CollMedicine
} % End of \newcommand{}

% ----------------------------------------------
%                    護理 Nursing
% ----------------------------------------------

% 護理學系 Department of Nursing
\newcommand{\SetDeptNursing}[0]
{
  \SetDeptName{護理學系}{Nursing}{Department of Nursing}
  \CollMedicine
} % End of \newcommand{}

% 醫學檢驗生物技術學系 Department of Medical Laboratory Science and Biotechnology
\newcommand{\SetDeptMT}[0]
{
  \SetDeptName{醫學檢驗生物技術學系}{MT}{Department of Medical Laboratory \\Science and Biotechnology}
  \CollMedicine
} % End of \newcommand{}

% 物理治療學系 Department of Physical Therapy
\newcommand{\SetDeptPT}[0]
{
  \SetDeptName{物理治療學系}{PT}{Department of Physical Therapy}
  \CollMedicine
} % End of \newcommand{}

% 職能治療學系 Department of Occupational Therapy
\newcommand{\SetDeptOT}[0]
{
  \SetDeptName{職能治療學系}{OT}{Department of Occupational Therapy}
  \CollMedicine
} % End of \newcommand{}

% 藥學系 School of Pharmacy
\newcommand{\SetDeptPharmacy}[0]
{
  \SetDeptName{藥學系}{Pharmacy}{School of Pharmacy}
  \CollMedicine
} % End of \newcommand{}

% 基礎醫學研究所 Institute of Basic Medical Sciences
\newcommand{\SetDeptBasicMed}[0]
{
  \SetDeptName{基礎醫學研究所}{BasicMed}{Institute of Basic Medical Sciences}
  \CollMedicine
} % End of \newcommand{}

% 行為醫學研究所 Institute of Behavioral Medicine
\newcommand{\SetDeptBehMed}[0]
{
  \SetDeptName{行為醫學研究所}{BehMed}{Institute of Behavioral Medicine}
  \CollMedicine
} % End of \newcommand{}

% 臨床藥學與藥物科技研究所 Institute of Clinical Pharmacy and Pharmaceutical Sciences
\newcommand{\SetDeptCLPARM}[0]
{
  \SetDeptName{臨床藥學與藥物科技研究所}{CLPARM}{Institute of Clinical Pharmacy \\ and Pharmaceutical Sciences}
  \CollMedicine
} % End of \newcommand{}

% 分子醫學研究所 Institute of Molecular Medicine
\newcommand{\SetDeptIMM}[0]
{
  \SetDeptName{分子醫學研究所}{IMM}{Institute of Molecular Medicine}
  \CollMedicine
} % End of \newcommand{}

% 口腔醫學研究所 Institute of Oral Medicine
\newcommand{\SetDeptIOM}[0]
{
  \SetDeptName{口腔醫學研究所}{IOM}{Institute of Oral Medicine}
  \CollMedicine
} % End of \newcommand{}

% 臨床醫學研究所 Institute of Clinical Medicine
\newcommand{\SetDeptICMMed}[0]
{
  \SetDeptName{臨床醫學研究所}{ICMMed}{Institute of Clinical Medicine}
  \CollMedicine
} % End of \newcommand{}

% 健康照護科學研究所 Institute of Allied Health Sciences
\newcommand{\SetDeptAlliedHealth}[0]
{
  \SetDeptName{健康照護科學研究所}{AlliedHealth}{Institute of Allied Health Sciences}
  \CollMedicine
} % End of \newcommand{}

% 老年學研究所 Institute of Gerontology
\newcommand{\SetDeptIOG}[0]
{
  \SetDeptName{老年學研究所}{IOG}{Institute of Gerontology}
  \CollMedicine
} % End of \newcommand{}

% ----------------------------------------------------------------------------


%
% This file is part of ncku-thesis-templete.
%
% ncku-thesis-templete is distributed in the hope that it will be useful,
% you can redistribute it and/or modify
% it under the terms of the Attribution-NonCommercial-ShareAlike
% 4.0 International.
%
% You should have received a copy of the
% Attribution-NonCommercial-ShareAlike 4.0 International
% along with ncku-thesis-templete.
%
% If not, see <http://creativecommons.org/licenses/by-nc-sa/4.0/legalcode.txt>.
%

% Some helper function about watermark

% ----------------------------------------------------------------------------

% 學校Logo watermark
\newcommand{\UseSchoolWatermark}[0]
{
  \AddToShipoutPicture{
  \put(0,0){
  \parbox[b][\paperheight]{\paperwidth}{
  \vfill
  \centering
  \includegraphics[]{./ncku/watersymbol.jpg}
  \vfill
  }}}
} % End of \newcommand{}

% 學校文字的watermark
\newcommand{\UseSchoolTextWatermark}[0]
{
  \AddToShipoutPicture*{
  \put(0,0){
  \parbox[b][\paperheight]{\paperwidth}{
  \vfill
  \centering
   \makebox(0,0){\rotatebox{45}{\textcolor[gray]{0.75}%
          {\fontsize{2.0cm}{2.0cm}\selectfont{\GetSchoolChiName}}}}
  \vfill
  }}}
} % End of \newcommand{}

\newcommand{\ClearWatermark}[0]{\ClearShipoutPicture}

% ----------------------------------------------------------------------------



% Language
%
% This file is part of ncku-thesis-templete.
%
% ncku-thesis-templete is distributed in the hope that it will be useful,
% you can redistribute it and/or modify
% it under the terms of the Attribution-NonCommercial-ShareAlike
% 4.0 International.
%
% You should have received a copy of the
% Attribution-NonCommercial-ShareAlike 4.0 International
% along with ncku-thesis-templete.
%
% If not, see <http://creativecommons.org/licenses/by-nc-sa/4.0/legalcode.txt>.
%

% Some helper function about language

% ----------------------------------------------------------------------------

% Flag
\newcommand\LangEnableChi{0}  % Default
\newcommand\LangEnableEng{0}  % Default
\newcommand\LangEnableDemo{0} % Default
\newcommand{\SetLangEnableChi}[0]{\renewcommand{\LangEnableChi}{1}}
\newcommand{\SetLangEnableEng}[0]{\renewcommand{\LangEnableEng}{1}}
\newcommand{\SetLangEnableDemo}[0]{\renewcommand{\LangEnableDemo}{1}}
\newcommand{\SetLangEnableAll}[0]
{
  \SetLangEnableChi
  \SetLangEnableEng
  \SetLangEnableDemo
} % End of \newcommand{}
\newcommand{\GetLangEnableChi}[0]{\LangEnableChi}
\newcommand{\GetLangEnableEng}[0]{\LangEnableEng}
\newcommand{\GetLangEnableDemo}[0]{\LangEnableDemo}

% APIs for user to do configure
\newcommand{\LangChi}[0]{\SetLangEnableChi}
\newcommand{\LangEng}[0]{\SetLangEnableEng}
\newcommand{\LangDemo}[0]{\SetLangEnableAll}

% ----------------------------------------------------------------------------

% Mode
\newcommand{\ChiMode}{\LangChi}
\newcommand{\EngMode}{\LangEng}
\newcommand{\DemoMode}{\LangDemo}

% ----------------------------------------------------------------------------


% ----------------------------------------------------------------------------



% --------------------------

% 論文有關資料

% This file is need to encoded in utf-8
%
% Choose or fill in some needed information for this thesis or dissertation
% 選擇或填入你的論文一些需要使用的資料

% ----------------------------------------------------------------------------

% --- 論文的編寫語言 / 用途 ---
% 只能選擇其中一個

  \ChiMode %:   編寫的為中英混合版, 有提供基本的所需的檔案
% \EngMode:  編寫的為全英文版, 有提供基本的所需的檔案
% \DemoMode: 完整教學 或 樣板測試用

% \DemoMode

%--------------------------------------------------

% --- 封面名字顯示方式 ---
% (如果你是使用\ChiMode, 則可無視這個設定)

% 如果你選擇是\EngMode,
% 預設在封面上只會顯示英文名字而已
% 把'%'拿掉以設定同時顯示中英文名字

%\CDBothName

% ----------------------------------------------------------------------------

% --- Title 論文題目 ---
% 填寫中文和(或)英文
% 如果題目內有必須以數學模式表示的符號，請用 \mbox{} 包住數學模式
% 如果覺得自動產生出來的題目斷行位置不適合, 可以手動加'\\'來強制斷行
% (圖書館說不管你是編寫中英混合或全英文版, 都必須同時存在中英題目)
%
% 有3種可使用, 可獨立使用, 但只有最後設定的一方有效
% \SetTitle{你的題目}{Your Title}   % 同時設定中英文題目
% \SetChiTitle{你的題目}            % 只設定中文題目
% \SetEngTitle{Your Title}         % 只設定英文題目
%
% e.g:
%
% \SetTitle %
% {國立成功大學碩博士用畢業論文XeLaTex模板} %
% {National Cheng Kung University (NCKU) Thesis/Dissertation Template in XeLaTex}
%
% or
%
% \SetChiTitle{國立成功大學碩博士用畢業論文XeLaTex模板}
% \SetEngTitle{National Cheng Kung University (NCKU)\\ Thesis/Dissertation Template in XeLaTex}

\SetTitle %
{社群組織之重要構成因素探討:成員分享行為及輔助分享資訊工具研究 \\ ---以g0v零時政府為例} %
{The Research of important constitution factor of Web community :Member sharing behavior and aid communication tool study\\ ---A g0v.com as an  example }


% ----------------------------------------------------------------------------

% --- Degree name 學位 ---
%
% 有2種可選擇, 但只有最後設定的一方有效
% \PhdDegree    % 博士學位
 \MasterDegree % 碩士學位

%\PhdDegree

% ----------------------------------------------------------------------------

% --- Your name 你的名字 ---
% 填寫你的中文和(或)英文

% 有3種可使用, 可獨立使用, 但只有最後設定的一方有效
% \SetMyName{你的名字}{Your name}   % 同時設定你的中英文名字
% \SetMyChiName{你的名字}           % 只設定你的中文名字
% \SetMyEngName{Your name}         % 只設定你的英文名字

\SetMyName{郭依如}{Lilian Kuo}

% ----------------------------------------------------------------------------

% --- Date 日期 ---

% 論文封面年月日因教務處仍無最後之確定版本, 所以目前仍以口試合格日為基準點.
% (若口試日期為 2014/10/01, 則封面日期需為2014/10, 2014/10/01或之後日期)

% --- 論文封面上的日期 ---
% 設定西元的年月, 會自動計算出民國的年份, 和英文的月份轉換
% 次序: {年份}{月份}
% \SetThesisDate{2014}{12}

\SetThesisDate{2015}{6}

%--------------------------------------------------

% --- 口試的日期 ---
% 設定西元的年月日, 會自動計算出民國的年份, 和英文的月份轉換
% 次序: {年份}{月份}{日}
% \SetOralDate{2014}{12}{31}

\SetOralDate{2015}{06}{20}

% ----------------------------------------------------------------------------

% --- 系所 Department or Institute ---
%
% 設定你的系所名字, e.g:
% \SetDeptMath 數學系
% \SetDeptCSIE 資訊工程學系

%\SetDeptEMBA %高階管理碩士在職專班

% ----------------------------------------------------------------------------

% --- 指導老師 Advisor(s) ---
% 在封面上預算了最多3位的空間
% 中文名字固定以 教授  為結尾
% 英文名字固定以 Prof. 為開頭

% 有3種可使用, 用來設定3位老師的名字
% \SetAdvisorNameX{老師的名字}{Professor's name} % 同時設定中英文名字
% \SetAdvisorChiNameX{老師的名字}                % 只設定中文名字
%\SetAdvisorChiNameX{王本正}          % 只設定中文名字
% \SetAdvisorEngNameX{Professor's name}         % 只設定英文名字
% (NameX為NameA, NameB, NameC)

% 使用\SetAdvisorNameA是必須的, 而如果你的指導教授有2或3位,
% 那只要增加\SetAdvisorNameB和\SetAdvisorNameC則可

\SetAdvisorNameA{王本正}{Kenny Wang}
%\SetAdvisorNameB{B}{B}
%\SetAdvisorNameC{C}{C}

% ----------------------------------------------------------------------------

% --- 口試証明文件 Oral presentation document ---
% 範例版本 或 使用檔案 只能選擇其中一方

% 顯示 / 使用 口試範例版本
%\DisplayOralTemplete
% --- 口試委員 Committee member(s) ---
% 口試委員數量 (至少2位, 最多8位, 預設為8位) 
% (Only work with \DisplayOralTemplete)
% 博士學位考試委員會置委員五人至九人
% 碩士學位考試委員會置委員三人至五人
% 口試委員人數含指導教授
%\SetCommitteeSize{5}

%--------------------------------------------------

% 使用口試圖片檔案
% 把你的圖片放在'context/oral'下
% 之後設定中英文版所對應是哪一個檔案
% (例子用的'oral-chi.pdf'和'oral-eng.pdf'已放在'context/oral'中)
%
%\DisplayOralImage                % 顯示圖片
%\SetOralImageChi{oral-chi.pdf}   % 中文口試檔案
%\SetOralImageEng{oral-eng.pdf}   % 英文口試檔案

% ----------------------------------------------------------------------------

% 關鍵字 Keyword
% 最多5個關鍵字

%\SetKeywords{Keyword A}{Keyword B}{Keyword C}{Keyword D}{Keyword E}

\SetKeywords{社群}{分享}{g0v}{非營利}{組織建構}

% ----------------------------------------------------------------------------

% 書脊 Spine

% 預設書脊上會使用英文題目
% 把'%'拿掉以設定改使用中文題目
% Default using the english title in Spine
% Remove '%' for change to use chinese title

\SpineTitleChi

% ----------------------------------------------------------------------------



% --------------------------

% 設定學校浮水印 Watermark
%\UseSchoolWatermark

% --------------------------

% 在 pdf 簡介欄裡填入相關資料
\FillInPDFData

% ------------------------------------------------

